\pdfoutput=1 \pdfcompresslevel=0 \pdfobjcompresslevel=0
\documentclass[11pt,a4paper]{report}
\year=2026 \month=8 \day=19 \time=600
\usepackage[margin=1in]{geometry}
\usepackage{amsmath,amssymb,amsthm,mathtools}
\usepackage[protrusion=true,expansion=false]{microtype}
\usepackage[colorlinks=true,hypertexnames=false]{hyperref}
\usepackage{xr}
\usepackage[capitalize,noabbrev]{cleveref}
\theoremstyle{plain}
\newtheorem{proposition}{Proposition}[chapter]
\newcommand{\vecn}[1]{\mathbf{#1}}
\newcommand{\pairs}[2]{\langle #1,#2\rangle}
\emergencystretch=3em
\begin{document}
\tableofcontents
\chapter{A chapter of ordinary matter, number 1}
\section{Prose that fills lines and breaks pages}
\label{sec:prose1-1}
Ordinary running text is what a typesetting engine spends most of its effort
upon, so the training document supplies a great deal of it. The words are
chosen for their length and their hyphenation, with incomprehensibility and
antidisestablishmentarianism thrown in so that the patterns are exercised.
Cross references such as \cref{sec:prose1-1} and citations of a page such
as \pageref{sec:prose1-1} make the auxiliary machinery run, and a
paragraph long enough to break over several lines gives the line breaker
something to weigh.
\begin{proposition}\label{prop:1-1}
For every $n\in\mathbb{N}$ the sum $\sum_{k=1}^{n}k$ equals
$\tfrac{1}{2}n(n+1)$.
\end{proposition}
\begin{proof}
Induction on $n$; the base case is $\pairs{1}{1}$ and the step adds $n+1$ to
both sides of the identity, which is what \cref{prop:1-1} asserts.
\end{proof}
\begin{equation}\label{eq:1-1}
  \int_{0}^{\infty}\frac{\vecn{x}^{s-1}}{e^{\vecn{x}}-1}\,d\vecn{x}
    = \Gamma(s)\zeta(s),\qquad \Re(s)>1 .
\end{equation}
\begin{align}
  \mathcal{L}[f](s) &= \int_{0}^{\infty} e^{-st} f(t)\,dt, \\
  \widehat{f}(\xi)  &= \int_{-\infty}^{\infty} f(x) e^{-2\pi i x\xi}\,dx .
\end{align}
\par\noindent\begin{center}
\begin{tabular}{lrrr}
\hline
Name & One & Two & Three \\
\hline
Alpha & 1 & 2 & 3 \\
Beta & 4 & 5 & 6 \\
Gamma & 7 & 8 & 9 \\
\hline
\end{tabular}
\end{center}
\begin{itemize}
  \item A first item with a formula $a^2+b^2=c^2$ in it.
  \item A second item that refers to \eqref{eq:1-1}.
  \item A third item, longer, so that the list breaks across a line and the
    paragraph builder has something to weigh when it chooses the break.
\end{itemize}
\subsection{A subsection}
Running text again, with \textit{italics}, \textbf{bold}, and
\texttt{typewriter}, and a formula $\alpha\beta\gamma\delta$ set inline so
that the math list machinery runs inside a paragraph. A displayed formula
follows:
\[ \lim_{n\to\infty}\left(1+\frac{1}{n}\right)^{n}=e . \]
And a little more prose to fill the page out, with an em-dash --- and
quotation marks ``like these'' and a footnote.\footnote{A footnote, which is
an insertion, so the page builder must split it.}
\subsection{A subsection}
Running text again, with \textit{italics}, \textbf{bold}, and
\texttt{typewriter}, and a formula $\alpha\beta\gamma\delta$ set inline so
that the math list machinery runs inside a paragraph. A displayed formula
follows:
\[ \lim_{n\to\infty}\left(1+\frac{1}{n}\right)^{n}=e . \]
And a little more prose to fill the page out, with an em-dash --- and
quotation marks ``like these'' and a footnote.\footnote{A footnote, which is
an insertion, so the page builder must split it.}
\subsection{A subsection}
Running text again, with \textit{italics}, \textbf{bold}, and
\texttt{typewriter}, and a formula $\alpha\beta\gamma\delta$ set inline so
that the math list machinery runs inside a paragraph. A displayed formula
follows:
\[ \lim_{n\to\infty}\left(1+\frac{1}{n}\right)^{n}=e . \]
And a little more prose to fill the page out, with an em-dash --- and
quotation marks ``like these'' and a footnote.\footnote{A footnote, which is
an insertion, so the page builder must split it.}
\section{Prose that fills lines and breaks pages}
\label{sec:prose1-2}
Ordinary running text is what a typesetting engine spends most of its effort
upon, so the training document supplies a great deal of it. The words are
chosen for their length and their hyphenation, with incomprehensibility and
antidisestablishmentarianism thrown in so that the patterns are exercised.
Cross references such as \cref{sec:prose1-2} and citations of a page such
as \pageref{sec:prose1-2} make the auxiliary machinery run, and a
paragraph long enough to break over several lines gives the line breaker
something to weigh.
\begin{proposition}\label{prop:1-2}
For every $n\in\mathbb{N}$ the sum $\sum_{k=1}^{n}k$ equals
$\tfrac{1}{2}n(n+1)$.
\end{proposition}
\begin{proof}
Induction on $n$; the base case is $\pairs{1}{1}$ and the step adds $n+1$ to
both sides of the identity, which is what \cref{prop:1-2} asserts.
\end{proof}
\begin{equation}\label{eq:1-2}
  \int_{0}^{\infty}\frac{\vecn{x}^{s-1}}{e^{\vecn{x}}-1}\,d\vecn{x}
    = \Gamma(s)\zeta(s),\qquad \Re(s)>1 .
\end{equation}
\begin{align}
  \mathcal{L}[f](s) &= \int_{0}^{\infty} e^{-st} f(t)\,dt, \\
  \widehat{f}(\xi)  &= \int_{-\infty}^{\infty} f(x) e^{-2\pi i x\xi}\,dx .
\end{align}
\par\noindent\begin{center}
\begin{tabular}{lrrr}
\hline
Name & One & Two & Three \\
\hline
Alpha & 1 & 2 & 3 \\
Beta & 4 & 5 & 6 \\
Gamma & 7 & 8 & 9 \\
\hline
\end{tabular}
\end{center}
\begin{itemize}
  \item A first item with a formula $a^2+b^2=c^2$ in it.
  \item A second item that refers to \eqref{eq:1-2}.
  \item A third item, longer, so that the list breaks across a line and the
    paragraph builder has something to weigh when it chooses the break.
\end{itemize}
\subsection{A subsection}
Running text again, with \textit{italics}, \textbf{bold}, and
\texttt{typewriter}, and a formula $\alpha\beta\gamma\delta$ set inline so
that the math list machinery runs inside a paragraph. A displayed formula
follows:
\[ \lim_{n\to\infty}\left(1+\frac{1}{n}\right)^{n}=e . \]
And a little more prose to fill the page out, with an em-dash --- and
quotation marks ``like these'' and a footnote.\footnote{A footnote, which is
an insertion, so the page builder must split it.}
\subsection{A subsection}
Running text again, with \textit{italics}, \textbf{bold}, and
\texttt{typewriter}, and a formula $\alpha\beta\gamma\delta$ set inline so
that the math list machinery runs inside a paragraph. A displayed formula
follows:
\[ \lim_{n\to\infty}\left(1+\frac{1}{n}\right)^{n}=e . \]
And a little more prose to fill the page out, with an em-dash --- and
quotation marks ``like these'' and a footnote.\footnote{A footnote, which is
an insertion, so the page builder must split it.}
\subsection{A subsection}
Running text again, with \textit{italics}, \textbf{bold}, and
\texttt{typewriter}, and a formula $\alpha\beta\gamma\delta$ set inline so
that the math list machinery runs inside a paragraph. A displayed formula
follows:
\[ \lim_{n\to\infty}\left(1+\frac{1}{n}\right)^{n}=e . \]
And a little more prose to fill the page out, with an em-dash --- and
quotation marks ``like these'' and a footnote.\footnote{A footnote, which is
an insertion, so the page builder must split it.}
\section{Prose that fills lines and breaks pages}
\label{sec:prose1-3}
Ordinary running text is what a typesetting engine spends most of its effort
upon, so the training document supplies a great deal of it. The words are
chosen for their length and their hyphenation, with incomprehensibility and
antidisestablishmentarianism thrown in so that the patterns are exercised.
Cross references such as \cref{sec:prose1-3} and citations of a page such
as \pageref{sec:prose1-3} make the auxiliary machinery run, and a
paragraph long enough to break over several lines gives the line breaker
something to weigh.
\begin{proposition}\label{prop:1-3}
For every $n\in\mathbb{N}$ the sum $\sum_{k=1}^{n}k$ equals
$\tfrac{1}{2}n(n+1)$.
\end{proposition}
\begin{proof}
Induction on $n$; the base case is $\pairs{1}{1}$ and the step adds $n+1$ to
both sides of the identity, which is what \cref{prop:1-3} asserts.
\end{proof}
\begin{equation}\label{eq:1-3}
  \int_{0}^{\infty}\frac{\vecn{x}^{s-1}}{e^{\vecn{x}}-1}\,d\vecn{x}
    = \Gamma(s)\zeta(s),\qquad \Re(s)>1 .
\end{equation}
\begin{align}
  \mathcal{L}[f](s) &= \int_{0}^{\infty} e^{-st} f(t)\,dt, \\
  \widehat{f}(\xi)  &= \int_{-\infty}^{\infty} f(x) e^{-2\pi i x\xi}\,dx .
\end{align}
\par\noindent\begin{center}
\begin{tabular}{lrrr}
\hline
Name & One & Two & Three \\
\hline
Alpha & 1 & 2 & 3 \\
Beta & 4 & 5 & 6 \\
Gamma & 7 & 8 & 9 \\
\hline
\end{tabular}
\end{center}
\begin{itemize}
  \item A first item with a formula $a^2+b^2=c^2$ in it.
  \item A second item that refers to \eqref{eq:1-3}.
  \item A third item, longer, so that the list breaks across a line and the
    paragraph builder has something to weigh when it chooses the break.
\end{itemize}
\subsection{A subsection}
Running text again, with \textit{italics}, \textbf{bold}, and
\texttt{typewriter}, and a formula $\alpha\beta\gamma\delta$ set inline so
that the math list machinery runs inside a paragraph. A displayed formula
follows:
\[ \lim_{n\to\infty}\left(1+\frac{1}{n}\right)^{n}=e . \]
And a little more prose to fill the page out, with an em-dash --- and
quotation marks ``like these'' and a footnote.\footnote{A footnote, which is
an insertion, so the page builder must split it.}
\subsection{A subsection}
Running text again, with \textit{italics}, \textbf{bold}, and
\texttt{typewriter}, and a formula $\alpha\beta\gamma\delta$ set inline so
that the math list machinery runs inside a paragraph. A displayed formula
follows:
\[ \lim_{n\to\infty}\left(1+\frac{1}{n}\right)^{n}=e . \]
And a little more prose to fill the page out, with an em-dash --- and
quotation marks ``like these'' and a footnote.\footnote{A footnote, which is
an insertion, so the page builder must split it.}
\subsection{A subsection}
Running text again, with \textit{italics}, \textbf{bold}, and
\texttt{typewriter}, and a formula $\alpha\beta\gamma\delta$ set inline so
that the math list machinery runs inside a paragraph. A displayed formula
follows:
\[ \lim_{n\to\infty}\left(1+\frac{1}{n}\right)^{n}=e . \]
And a little more prose to fill the page out, with an em-dash --- and
quotation marks ``like these'' and a footnote.\footnote{A footnote, which is
an insertion, so the page builder must split it.}
\section{Prose that fills lines and breaks pages}
\label{sec:prose1-4}
Ordinary running text is what a typesetting engine spends most of its effort
upon, so the training document supplies a great deal of it. The words are
chosen for their length and their hyphenation, with incomprehensibility and
antidisestablishmentarianism thrown in so that the patterns are exercised.
Cross references such as \cref{sec:prose1-4} and citations of a page such
as \pageref{sec:prose1-4} make the auxiliary machinery run, and a
paragraph long enough to break over several lines gives the line breaker
something to weigh.
\begin{proposition}\label{prop:1-4}
For every $n\in\mathbb{N}$ the sum $\sum_{k=1}^{n}k$ equals
$\tfrac{1}{2}n(n+1)$.
\end{proposition}
\begin{proof}
Induction on $n$; the base case is $\pairs{1}{1}$ and the step adds $n+1$ to
both sides of the identity, which is what \cref{prop:1-4} asserts.
\end{proof}
\begin{equation}\label{eq:1-4}
  \int_{0}^{\infty}\frac{\vecn{x}^{s-1}}{e^{\vecn{x}}-1}\,d\vecn{x}
    = \Gamma(s)\zeta(s),\qquad \Re(s)>1 .
\end{equation}
\begin{align}
  \mathcal{L}[f](s) &= \int_{0}^{\infty} e^{-st} f(t)\,dt, \\
  \widehat{f}(\xi)  &= \int_{-\infty}^{\infty} f(x) e^{-2\pi i x\xi}\,dx .
\end{align}
\par\noindent\begin{center}
\begin{tabular}{lrrr}
\hline
Name & One & Two & Three \\
\hline
Alpha & 1 & 2 & 3 \\
Beta & 4 & 5 & 6 \\
Gamma & 7 & 8 & 9 \\
\hline
\end{tabular}
\end{center}
\begin{itemize}
  \item A first item with a formula $a^2+b^2=c^2$ in it.
  \item A second item that refers to \eqref{eq:1-4}.
  \item A third item, longer, so that the list breaks across a line and the
    paragraph builder has something to weigh when it chooses the break.
\end{itemize}
\subsection{A subsection}
Running text again, with \textit{italics}, \textbf{bold}, and
\texttt{typewriter}, and a formula $\alpha\beta\gamma\delta$ set inline so
that the math list machinery runs inside a paragraph. A displayed formula
follows:
\[ \lim_{n\to\infty}\left(1+\frac{1}{n}\right)^{n}=e . \]
And a little more prose to fill the page out, with an em-dash --- and
quotation marks ``like these'' and a footnote.\footnote{A footnote, which is
an insertion, so the page builder must split it.}
\subsection{A subsection}
Running text again, with \textit{italics}, \textbf{bold}, and
\texttt{typewriter}, and a formula $\alpha\beta\gamma\delta$ set inline so
that the math list machinery runs inside a paragraph. A displayed formula
follows:
\[ \lim_{n\to\infty}\left(1+\frac{1}{n}\right)^{n}=e . \]
And a little more prose to fill the page out, with an em-dash --- and
quotation marks ``like these'' and a footnote.\footnote{A footnote, which is
an insertion, so the page builder must split it.}
\subsection{A subsection}
Running text again, with \textit{italics}, \textbf{bold}, and
\texttt{typewriter}, and a formula $\alpha\beta\gamma\delta$ set inline so
that the math list machinery runs inside a paragraph. A displayed formula
follows:
\[ \lim_{n\to\infty}\left(1+\frac{1}{n}\right)^{n}=e . \]
And a little more prose to fill the page out, with an em-dash --- and
quotation marks ``like these'' and a footnote.\footnote{A footnote, which is
an insertion, so the page builder must split it.}
\chapter{A chapter of ordinary matter, number 2}
\section{Prose that fills lines and breaks pages}
\label{sec:prose2-1}
Ordinary running text is what a typesetting engine spends most of its effort
upon, so the training document supplies a great deal of it. The words are
chosen for their length and their hyphenation, with incomprehensibility and
antidisestablishmentarianism thrown in so that the patterns are exercised.
Cross references such as \cref{sec:prose2-1} and citations of a page such
as \pageref{sec:prose2-1} make the auxiliary machinery run, and a
paragraph long enough to break over several lines gives the line breaker
something to weigh.
\begin{proposition}\label{prop:2-1}
For every $n\in\mathbb{N}$ the sum $\sum_{k=1}^{n}k$ equals
$\tfrac{1}{2}n(n+1)$.
\end{proposition}
\begin{proof}
Induction on $n$; the base case is $\pairs{1}{1}$ and the step adds $n+1$ to
both sides of the identity, which is what \cref{prop:2-1} asserts.
\end{proof}
\begin{equation}\label{eq:2-1}
  \int_{0}^{\infty}\frac{\vecn{x}^{s-1}}{e^{\vecn{x}}-1}\,d\vecn{x}
    = \Gamma(s)\zeta(s),\qquad \Re(s)>1 .
\end{equation}
\begin{align}
  \mathcal{L}[f](s) &= \int_{0}^{\infty} e^{-st} f(t)\,dt, \\
  \widehat{f}(\xi)  &= \int_{-\infty}^{\infty} f(x) e^{-2\pi i x\xi}\,dx .
\end{align}
\par\noindent\begin{center}
\begin{tabular}{lrrr}
\hline
Name & One & Two & Three \\
\hline
Alpha & 1 & 2 & 3 \\
Beta & 4 & 5 & 6 \\
Gamma & 7 & 8 & 9 \\
\hline
\end{tabular}
\end{center}
\begin{itemize}
  \item A first item with a formula $a^2+b^2=c^2$ in it.
  \item A second item that refers to \eqref{eq:2-1}.
  \item A third item, longer, so that the list breaks across a line and the
    paragraph builder has something to weigh when it chooses the break.
\end{itemize}
\subsection{A subsection}
Running text again, with \textit{italics}, \textbf{bold}, and
\texttt{typewriter}, and a formula $\alpha\beta\gamma\delta$ set inline so
that the math list machinery runs inside a paragraph. A displayed formula
follows:
\[ \lim_{n\to\infty}\left(1+\frac{1}{n}\right)^{n}=e . \]
And a little more prose to fill the page out, with an em-dash --- and
quotation marks ``like these'' and a footnote.\footnote{A footnote, which is
an insertion, so the page builder must split it.}
\subsection{A subsection}
Running text again, with \textit{italics}, \textbf{bold}, and
\texttt{typewriter}, and a formula $\alpha\beta\gamma\delta$ set inline so
that the math list machinery runs inside a paragraph. A displayed formula
follows:
\[ \lim_{n\to\infty}\left(1+\frac{1}{n}\right)^{n}=e . \]
And a little more prose to fill the page out, with an em-dash --- and
quotation marks ``like these'' and a footnote.\footnote{A footnote, which is
an insertion, so the page builder must split it.}
\subsection{A subsection}
Running text again, with \textit{italics}, \textbf{bold}, and
\texttt{typewriter}, and a formula $\alpha\beta\gamma\delta$ set inline so
that the math list machinery runs inside a paragraph. A displayed formula
follows:
\[ \lim_{n\to\infty}\left(1+\frac{1}{n}\right)^{n}=e . \]
And a little more prose to fill the page out, with an em-dash --- and
quotation marks ``like these'' and a footnote.\footnote{A footnote, which is
an insertion, so the page builder must split it.}
\section{Prose that fills lines and breaks pages}
\label{sec:prose2-2}
Ordinary running text is what a typesetting engine spends most of its effort
upon, so the training document supplies a great deal of it. The words are
chosen for their length and their hyphenation, with incomprehensibility and
antidisestablishmentarianism thrown in so that the patterns are exercised.
Cross references such as \cref{sec:prose2-2} and citations of a page such
as \pageref{sec:prose2-2} make the auxiliary machinery run, and a
paragraph long enough to break over several lines gives the line breaker
something to weigh.
\begin{proposition}\label{prop:2-2}
For every $n\in\mathbb{N}$ the sum $\sum_{k=1}^{n}k$ equals
$\tfrac{1}{2}n(n+1)$.
\end{proposition}
\begin{proof}
Induction on $n$; the base case is $\pairs{1}{1}$ and the step adds $n+1$ to
both sides of the identity, which is what \cref{prop:2-2} asserts.
\end{proof}
\begin{equation}\label{eq:2-2}
  \int_{0}^{\infty}\frac{\vecn{x}^{s-1}}{e^{\vecn{x}}-1}\,d\vecn{x}
    = \Gamma(s)\zeta(s),\qquad \Re(s)>1 .
\end{equation}
\begin{align}
  \mathcal{L}[f](s) &= \int_{0}^{\infty} e^{-st} f(t)\,dt, \\
  \widehat{f}(\xi)  &= \int_{-\infty}^{\infty} f(x) e^{-2\pi i x\xi}\,dx .
\end{align}
\par\noindent\begin{center}
\begin{tabular}{lrrr}
\hline
Name & One & Two & Three \\
\hline
Alpha & 1 & 2 & 3 \\
Beta & 4 & 5 & 6 \\
Gamma & 7 & 8 & 9 \\
\hline
\end{tabular}
\end{center}
\begin{itemize}
  \item A first item with a formula $a^2+b^2=c^2$ in it.
  \item A second item that refers to \eqref{eq:2-2}.
  \item A third item, longer, so that the list breaks across a line and the
    paragraph builder has something to weigh when it chooses the break.
\end{itemize}
\subsection{A subsection}
Running text again, with \textit{italics}, \textbf{bold}, and
\texttt{typewriter}, and a formula $\alpha\beta\gamma\delta$ set inline so
that the math list machinery runs inside a paragraph. A displayed formula
follows:
\[ \lim_{n\to\infty}\left(1+\frac{1}{n}\right)^{n}=e . \]
And a little more prose to fill the page out, with an em-dash --- and
quotation marks ``like these'' and a footnote.\footnote{A footnote, which is
an insertion, so the page builder must split it.}
\subsection{A subsection}
Running text again, with \textit{italics}, \textbf{bold}, and
\texttt{typewriter}, and a formula $\alpha\beta\gamma\delta$ set inline so
that the math list machinery runs inside a paragraph. A displayed formula
follows:
\[ \lim_{n\to\infty}\left(1+\frac{1}{n}\right)^{n}=e . \]
And a little more prose to fill the page out, with an em-dash --- and
quotation marks ``like these'' and a footnote.\footnote{A footnote, which is
an insertion, so the page builder must split it.}
\subsection{A subsection}
Running text again, with \textit{italics}, \textbf{bold}, and
\texttt{typewriter}, and a formula $\alpha\beta\gamma\delta$ set inline so
that the math list machinery runs inside a paragraph. A displayed formula
follows:
\[ \lim_{n\to\infty}\left(1+\frac{1}{n}\right)^{n}=e . \]
And a little more prose to fill the page out, with an em-dash --- and
quotation marks ``like these'' and a footnote.\footnote{A footnote, which is
an insertion, so the page builder must split it.}
\section{Prose that fills lines and breaks pages}
\label{sec:prose2-3}
Ordinary running text is what a typesetting engine spends most of its effort
upon, so the training document supplies a great deal of it. The words are
chosen for their length and their hyphenation, with incomprehensibility and
antidisestablishmentarianism thrown in so that the patterns are exercised.
Cross references such as \cref{sec:prose2-3} and citations of a page such
as \pageref{sec:prose2-3} make the auxiliary machinery run, and a
paragraph long enough to break over several lines gives the line breaker
something to weigh.
\begin{proposition}\label{prop:2-3}
For every $n\in\mathbb{N}$ the sum $\sum_{k=1}^{n}k$ equals
$\tfrac{1}{2}n(n+1)$.
\end{proposition}
\begin{proof}
Induction on $n$; the base case is $\pairs{1}{1}$ and the step adds $n+1$ to
both sides of the identity, which is what \cref{prop:2-3} asserts.
\end{proof}
\begin{equation}\label{eq:2-3}
  \int_{0}^{\infty}\frac{\vecn{x}^{s-1}}{e^{\vecn{x}}-1}\,d\vecn{x}
    = \Gamma(s)\zeta(s),\qquad \Re(s)>1 .
\end{equation}
\begin{align}
  \mathcal{L}[f](s) &= \int_{0}^{\infty} e^{-st} f(t)\,dt, \\
  \widehat{f}(\xi)  &= \int_{-\infty}^{\infty} f(x) e^{-2\pi i x\xi}\,dx .
\end{align}
\par\noindent\begin{center}
\begin{tabular}{lrrr}
\hline
Name & One & Two & Three \\
\hline
Alpha & 1 & 2 & 3 \\
Beta & 4 & 5 & 6 \\
Gamma & 7 & 8 & 9 \\
\hline
\end{tabular}
\end{center}
\begin{itemize}
  \item A first item with a formula $a^2+b^2=c^2$ in it.
  \item A second item that refers to \eqref{eq:2-3}.
  \item A third item, longer, so that the list breaks across a line and the
    paragraph builder has something to weigh when it chooses the break.
\end{itemize}
\subsection{A subsection}
Running text again, with \textit{italics}, \textbf{bold}, and
\texttt{typewriter}, and a formula $\alpha\beta\gamma\delta$ set inline so
that the math list machinery runs inside a paragraph. A displayed formula
follows:
\[ \lim_{n\to\infty}\left(1+\frac{1}{n}\right)^{n}=e . \]
And a little more prose to fill the page out, with an em-dash --- and
quotation marks ``like these'' and a footnote.\footnote{A footnote, which is
an insertion, so the page builder must split it.}
\subsection{A subsection}
Running text again, with \textit{italics}, \textbf{bold}, and
\texttt{typewriter}, and a formula $\alpha\beta\gamma\delta$ set inline so
that the math list machinery runs inside a paragraph. A displayed formula
follows:
\[ \lim_{n\to\infty}\left(1+\frac{1}{n}\right)^{n}=e . \]
And a little more prose to fill the page out, with an em-dash --- and
quotation marks ``like these'' and a footnote.\footnote{A footnote, which is
an insertion, so the page builder must split it.}
\subsection{A subsection}
Running text again, with \textit{italics}, \textbf{bold}, and
\texttt{typewriter}, and a formula $\alpha\beta\gamma\delta$ set inline so
that the math list machinery runs inside a paragraph. A displayed formula
follows:
\[ \lim_{n\to\infty}\left(1+\frac{1}{n}\right)^{n}=e . \]
And a little more prose to fill the page out, with an em-dash --- and
quotation marks ``like these'' and a footnote.\footnote{A footnote, which is
an insertion, so the page builder must split it.}
\section{Prose that fills lines and breaks pages}
\label{sec:prose2-4}
Ordinary running text is what a typesetting engine spends most of its effort
upon, so the training document supplies a great deal of it. The words are
chosen for their length and their hyphenation, with incomprehensibility and
antidisestablishmentarianism thrown in so that the patterns are exercised.
Cross references such as \cref{sec:prose2-4} and citations of a page such
as \pageref{sec:prose2-4} make the auxiliary machinery run, and a
paragraph long enough to break over several lines gives the line breaker
something to weigh.
\begin{proposition}\label{prop:2-4}
For every $n\in\mathbb{N}$ the sum $\sum_{k=1}^{n}k$ equals
$\tfrac{1}{2}n(n+1)$.
\end{proposition}
\begin{proof}
Induction on $n$; the base case is $\pairs{1}{1}$ and the step adds $n+1$ to
both sides of the identity, which is what \cref{prop:2-4} asserts.
\end{proof}
\begin{equation}\label{eq:2-4}
  \int_{0}^{\infty}\frac{\vecn{x}^{s-1}}{e^{\vecn{x}}-1}\,d\vecn{x}
    = \Gamma(s)\zeta(s),\qquad \Re(s)>1 .
\end{equation}
\begin{align}
  \mathcal{L}[f](s) &= \int_{0}^{\infty} e^{-st} f(t)\,dt, \\
  \widehat{f}(\xi)  &= \int_{-\infty}^{\infty} f(x) e^{-2\pi i x\xi}\,dx .
\end{align}
\par\noindent\begin{center}
\begin{tabular}{lrrr}
\hline
Name & One & Two & Three \\
\hline
Alpha & 1 & 2 & 3 \\
Beta & 4 & 5 & 6 \\
Gamma & 7 & 8 & 9 \\
\hline
\end{tabular}
\end{center}
\begin{itemize}
  \item A first item with a formula $a^2+b^2=c^2$ in it.
  \item A second item that refers to \eqref{eq:2-4}.
  \item A third item, longer, so that the list breaks across a line and the
    paragraph builder has something to weigh when it chooses the break.
\end{itemize}
\subsection{A subsection}
Running text again, with \textit{italics}, \textbf{bold}, and
\texttt{typewriter}, and a formula $\alpha\beta\gamma\delta$ set inline so
that the math list machinery runs inside a paragraph. A displayed formula
follows:
\[ \lim_{n\to\infty}\left(1+\frac{1}{n}\right)^{n}=e . \]
And a little more prose to fill the page out, with an em-dash --- and
quotation marks ``like these'' and a footnote.\footnote{A footnote, which is
an insertion, so the page builder must split it.}
\subsection{A subsection}
Running text again, with \textit{italics}, \textbf{bold}, and
\texttt{typewriter}, and a formula $\alpha\beta\gamma\delta$ set inline so
that the math list machinery runs inside a paragraph. A displayed formula
follows:
\[ \lim_{n\to\infty}\left(1+\frac{1}{n}\right)^{n}=e . \]
And a little more prose to fill the page out, with an em-dash --- and
quotation marks ``like these'' and a footnote.\footnote{A footnote, which is
an insertion, so the page builder must split it.}
\subsection{A subsection}
Running text again, with \textit{italics}, \textbf{bold}, and
\texttt{typewriter}, and a formula $\alpha\beta\gamma\delta$ set inline so
that the math list machinery runs inside a paragraph. A displayed formula
follows:
\[ \lim_{n\to\infty}\left(1+\frac{1}{n}\right)^{n}=e . \]
And a little more prose to fill the page out, with an em-dash --- and
quotation marks ``like these'' and a footnote.\footnote{A footnote, which is
an insertion, so the page builder must split it.}
\chapter{A chapter of ordinary matter, number 3}
\section{Prose that fills lines and breaks pages}
\label{sec:prose3-1}
Ordinary running text is what a typesetting engine spends most of its effort
upon, so the training document supplies a great deal of it. The words are
chosen for their length and their hyphenation, with incomprehensibility and
antidisestablishmentarianism thrown in so that the patterns are exercised.
Cross references such as \cref{sec:prose3-1} and citations of a page such
as \pageref{sec:prose3-1} make the auxiliary machinery run, and a
paragraph long enough to break over several lines gives the line breaker
something to weigh.
\begin{proposition}\label{prop:3-1}
For every $n\in\mathbb{N}$ the sum $\sum_{k=1}^{n}k$ equals
$\tfrac{1}{2}n(n+1)$.
\end{proposition}
\begin{proof}
Induction on $n$; the base case is $\pairs{1}{1}$ and the step adds $n+1$ to
both sides of the identity, which is what \cref{prop:3-1} asserts.
\end{proof}
\begin{equation}\label{eq:3-1}
  \int_{0}^{\infty}\frac{\vecn{x}^{s-1}}{e^{\vecn{x}}-1}\,d\vecn{x}
    = \Gamma(s)\zeta(s),\qquad \Re(s)>1 .
\end{equation}
\begin{align}
  \mathcal{L}[f](s) &= \int_{0}^{\infty} e^{-st} f(t)\,dt, \\
  \widehat{f}(\xi)  &= \int_{-\infty}^{\infty} f(x) e^{-2\pi i x\xi}\,dx .
\end{align}
\par\noindent\begin{center}
\begin{tabular}{lrrr}
\hline
Name & One & Two & Three \\
\hline
Alpha & 1 & 2 & 3 \\
Beta & 4 & 5 & 6 \\
Gamma & 7 & 8 & 9 \\
\hline
\end{tabular}
\end{center}
\begin{itemize}
  \item A first item with a formula $a^2+b^2=c^2$ in it.
  \item A second item that refers to \eqref{eq:3-1}.
  \item A third item, longer, so that the list breaks across a line and the
    paragraph builder has something to weigh when it chooses the break.
\end{itemize}
\subsection{A subsection}
Running text again, with \textit{italics}, \textbf{bold}, and
\texttt{typewriter}, and a formula $\alpha\beta\gamma\delta$ set inline so
that the math list machinery runs inside a paragraph. A displayed formula
follows:
\[ \lim_{n\to\infty}\left(1+\frac{1}{n}\right)^{n}=e . \]
And a little more prose to fill the page out, with an em-dash --- and
quotation marks ``like these'' and a footnote.\footnote{A footnote, which is
an insertion, so the page builder must split it.}
\subsection{A subsection}
Running text again, with \textit{italics}, \textbf{bold}, and
\texttt{typewriter}, and a formula $\alpha\beta\gamma\delta$ set inline so
that the math list machinery runs inside a paragraph. A displayed formula
follows:
\[ \lim_{n\to\infty}\left(1+\frac{1}{n}\right)^{n}=e . \]
And a little more prose to fill the page out, with an em-dash --- and
quotation marks ``like these'' and a footnote.\footnote{A footnote, which is
an insertion, so the page builder must split it.}
\subsection{A subsection}
Running text again, with \textit{italics}, \textbf{bold}, and
\texttt{typewriter}, and a formula $\alpha\beta\gamma\delta$ set inline so
that the math list machinery runs inside a paragraph. A displayed formula
follows:
\[ \lim_{n\to\infty}\left(1+\frac{1}{n}\right)^{n}=e . \]
And a little more prose to fill the page out, with an em-dash --- and
quotation marks ``like these'' and a footnote.\footnote{A footnote, which is
an insertion, so the page builder must split it.}
\section{Prose that fills lines and breaks pages}
\label{sec:prose3-2}
Ordinary running text is what a typesetting engine spends most of its effort
upon, so the training document supplies a great deal of it. The words are
chosen for their length and their hyphenation, with incomprehensibility and
antidisestablishmentarianism thrown in so that the patterns are exercised.
Cross references such as \cref{sec:prose3-2} and citations of a page such
as \pageref{sec:prose3-2} make the auxiliary machinery run, and a
paragraph long enough to break over several lines gives the line breaker
something to weigh.
\begin{proposition}\label{prop:3-2}
For every $n\in\mathbb{N}$ the sum $\sum_{k=1}^{n}k$ equals
$\tfrac{1}{2}n(n+1)$.
\end{proposition}
\begin{proof}
Induction on $n$; the base case is $\pairs{1}{1}$ and the step adds $n+1$ to
both sides of the identity, which is what \cref{prop:3-2} asserts.
\end{proof}
\begin{equation}\label{eq:3-2}
  \int_{0}^{\infty}\frac{\vecn{x}^{s-1}}{e^{\vecn{x}}-1}\,d\vecn{x}
    = \Gamma(s)\zeta(s),\qquad \Re(s)>1 .
\end{equation}
\begin{align}
  \mathcal{L}[f](s) &= \int_{0}^{\infty} e^{-st} f(t)\,dt, \\
  \widehat{f}(\xi)  &= \int_{-\infty}^{\infty} f(x) e^{-2\pi i x\xi}\,dx .
\end{align}
\par\noindent\begin{center}
\begin{tabular}{lrrr}
\hline
Name & One & Two & Three \\
\hline
Alpha & 1 & 2 & 3 \\
Beta & 4 & 5 & 6 \\
Gamma & 7 & 8 & 9 \\
\hline
\end{tabular}
\end{center}
\begin{itemize}
  \item A first item with a formula $a^2+b^2=c^2$ in it.
  \item A second item that refers to \eqref{eq:3-2}.
  \item A third item, longer, so that the list breaks across a line and the
    paragraph builder has something to weigh when it chooses the break.
\end{itemize}
\subsection{A subsection}
Running text again, with \textit{italics}, \textbf{bold}, and
\texttt{typewriter}, and a formula $\alpha\beta\gamma\delta$ set inline so
that the math list machinery runs inside a paragraph. A displayed formula
follows:
\[ \lim_{n\to\infty}\left(1+\frac{1}{n}\right)^{n}=e . \]
And a little more prose to fill the page out, with an em-dash --- and
quotation marks ``like these'' and a footnote.\footnote{A footnote, which is
an insertion, so the page builder must split it.}
\subsection{A subsection}
Running text again, with \textit{italics}, \textbf{bold}, and
\texttt{typewriter}, and a formula $\alpha\beta\gamma\delta$ set inline so
that the math list machinery runs inside a paragraph. A displayed formula
follows:
\[ \lim_{n\to\infty}\left(1+\frac{1}{n}\right)^{n}=e . \]
And a little more prose to fill the page out, with an em-dash --- and
quotation marks ``like these'' and a footnote.\footnote{A footnote, which is
an insertion, so the page builder must split it.}
\subsection{A subsection}
Running text again, with \textit{italics}, \textbf{bold}, and
\texttt{typewriter}, and a formula $\alpha\beta\gamma\delta$ set inline so
that the math list machinery runs inside a paragraph. A displayed formula
follows:
\[ \lim_{n\to\infty}\left(1+\frac{1}{n}\right)^{n}=e . \]
And a little more prose to fill the page out, with an em-dash --- and
quotation marks ``like these'' and a footnote.\footnote{A footnote, which is
an insertion, so the page builder must split it.}
\section{Prose that fills lines and breaks pages}
\label{sec:prose3-3}
Ordinary running text is what a typesetting engine spends most of its effort
upon, so the training document supplies a great deal of it. The words are
chosen for their length and their hyphenation, with incomprehensibility and
antidisestablishmentarianism thrown in so that the patterns are exercised.
Cross references such as \cref{sec:prose3-3} and citations of a page such
as \pageref{sec:prose3-3} make the auxiliary machinery run, and a
paragraph long enough to break over several lines gives the line breaker
something to weigh.
\begin{proposition}\label{prop:3-3}
For every $n\in\mathbb{N}$ the sum $\sum_{k=1}^{n}k$ equals
$\tfrac{1}{2}n(n+1)$.
\end{proposition}
\begin{proof}
Induction on $n$; the base case is $\pairs{1}{1}$ and the step adds $n+1$ to
both sides of the identity, which is what \cref{prop:3-3} asserts.
\end{proof}
\begin{equation}\label{eq:3-3}
  \int_{0}^{\infty}\frac{\vecn{x}^{s-1}}{e^{\vecn{x}}-1}\,d\vecn{x}
    = \Gamma(s)\zeta(s),\qquad \Re(s)>1 .
\end{equation}
\begin{align}
  \mathcal{L}[f](s) &= \int_{0}^{\infty} e^{-st} f(t)\,dt, \\
  \widehat{f}(\xi)  &= \int_{-\infty}^{\infty} f(x) e^{-2\pi i x\xi}\,dx .
\end{align}
\par\noindent\begin{center}
\begin{tabular}{lrrr}
\hline
Name & One & Two & Three \\
\hline
Alpha & 1 & 2 & 3 \\
Beta & 4 & 5 & 6 \\
Gamma & 7 & 8 & 9 \\
\hline
\end{tabular}
\end{center}
\begin{itemize}
  \item A first item with a formula $a^2+b^2=c^2$ in it.
  \item A second item that refers to \eqref{eq:3-3}.
  \item A third item, longer, so that the list breaks across a line and the
    paragraph builder has something to weigh when it chooses the break.
\end{itemize}
\subsection{A subsection}
Running text again, with \textit{italics}, \textbf{bold}, and
\texttt{typewriter}, and a formula $\alpha\beta\gamma\delta$ set inline so
that the math list machinery runs inside a paragraph. A displayed formula
follows:
\[ \lim_{n\to\infty}\left(1+\frac{1}{n}\right)^{n}=e . \]
And a little more prose to fill the page out, with an em-dash --- and
quotation marks ``like these'' and a footnote.\footnote{A footnote, which is
an insertion, so the page builder must split it.}
\subsection{A subsection}
Running text again, with \textit{italics}, \textbf{bold}, and
\texttt{typewriter}, and a formula $\alpha\beta\gamma\delta$ set inline so
that the math list machinery runs inside a paragraph. A displayed formula
follows:
\[ \lim_{n\to\infty}\left(1+\frac{1}{n}\right)^{n}=e . \]
And a little more prose to fill the page out, with an em-dash --- and
quotation marks ``like these'' and a footnote.\footnote{A footnote, which is
an insertion, so the page builder must split it.}
\subsection{A subsection}
Running text again, with \textit{italics}, \textbf{bold}, and
\texttt{typewriter}, and a formula $\alpha\beta\gamma\delta$ set inline so
that the math list machinery runs inside a paragraph. A displayed formula
follows:
\[ \lim_{n\to\infty}\left(1+\frac{1}{n}\right)^{n}=e . \]
And a little more prose to fill the page out, with an em-dash --- and
quotation marks ``like these'' and a footnote.\footnote{A footnote, which is
an insertion, so the page builder must split it.}
\section{Prose that fills lines and breaks pages}
\label{sec:prose3-4}
Ordinary running text is what a typesetting engine spends most of its effort
upon, so the training document supplies a great deal of it. The words are
chosen for their length and their hyphenation, with incomprehensibility and
antidisestablishmentarianism thrown in so that the patterns are exercised.
Cross references such as \cref{sec:prose3-4} and citations of a page such
as \pageref{sec:prose3-4} make the auxiliary machinery run, and a
paragraph long enough to break over several lines gives the line breaker
something to weigh.
\begin{proposition}\label{prop:3-4}
For every $n\in\mathbb{N}$ the sum $\sum_{k=1}^{n}k$ equals
$\tfrac{1}{2}n(n+1)$.
\end{proposition}
\begin{proof}
Induction on $n$; the base case is $\pairs{1}{1}$ and the step adds $n+1$ to
both sides of the identity, which is what \cref{prop:3-4} asserts.
\end{proof}
\begin{equation}\label{eq:3-4}
  \int_{0}^{\infty}\frac{\vecn{x}^{s-1}}{e^{\vecn{x}}-1}\,d\vecn{x}
    = \Gamma(s)\zeta(s),\qquad \Re(s)>1 .
\end{equation}
\begin{align}
  \mathcal{L}[f](s) &= \int_{0}^{\infty} e^{-st} f(t)\,dt, \\
  \widehat{f}(\xi)  &= \int_{-\infty}^{\infty} f(x) e^{-2\pi i x\xi}\,dx .
\end{align}
\par\noindent\begin{center}
\begin{tabular}{lrrr}
\hline
Name & One & Two & Three \\
\hline
Alpha & 1 & 2 & 3 \\
Beta & 4 & 5 & 6 \\
Gamma & 7 & 8 & 9 \\
\hline
\end{tabular}
\end{center}
\begin{itemize}
  \item A first item with a formula $a^2+b^2=c^2$ in it.
  \item A second item that refers to \eqref{eq:3-4}.
  \item A third item, longer, so that the list breaks across a line and the
    paragraph builder has something to weigh when it chooses the break.
\end{itemize}
\subsection{A subsection}
Running text again, with \textit{italics}, \textbf{bold}, and
\texttt{typewriter}, and a formula $\alpha\beta\gamma\delta$ set inline so
that the math list machinery runs inside a paragraph. A displayed formula
follows:
\[ \lim_{n\to\infty}\left(1+\frac{1}{n}\right)^{n}=e . \]
And a little more prose to fill the page out, with an em-dash --- and
quotation marks ``like these'' and a footnote.\footnote{A footnote, which is
an insertion, so the page builder must split it.}
\subsection{A subsection}
Running text again, with \textit{italics}, \textbf{bold}, and
\texttt{typewriter}, and a formula $\alpha\beta\gamma\delta$ set inline so
that the math list machinery runs inside a paragraph. A displayed formula
follows:
\[ \lim_{n\to\infty}\left(1+\frac{1}{n}\right)^{n}=e . \]
And a little more prose to fill the page out, with an em-dash --- and
quotation marks ``like these'' and a footnote.\footnote{A footnote, which is
an insertion, so the page builder must split it.}
\subsection{A subsection}
Running text again, with \textit{italics}, \textbf{bold}, and
\texttt{typewriter}, and a formula $\alpha\beta\gamma\delta$ set inline so
that the math list machinery runs inside a paragraph. A displayed formula
follows:
\[ \lim_{n\to\infty}\left(1+\frac{1}{n}\right)^{n}=e . \]
And a little more prose to fill the page out, with an em-dash --- and
quotation marks ``like these'' and a footnote.\footnote{A footnote, which is
an insertion, so the page builder must split it.}
\chapter{A chapter of ordinary matter, number 4}
\section{Prose that fills lines and breaks pages}
\label{sec:prose4-1}
Ordinary running text is what a typesetting engine spends most of its effort
upon, so the training document supplies a great deal of it. The words are
chosen for their length and their hyphenation, with incomprehensibility and
antidisestablishmentarianism thrown in so that the patterns are exercised.
Cross references such as \cref{sec:prose4-1} and citations of a page such
as \pageref{sec:prose4-1} make the auxiliary machinery run, and a
paragraph long enough to break over several lines gives the line breaker
something to weigh.
\begin{proposition}\label{prop:4-1}
For every $n\in\mathbb{N}$ the sum $\sum_{k=1}^{n}k$ equals
$\tfrac{1}{2}n(n+1)$.
\end{proposition}
\begin{proof}
Induction on $n$; the base case is $\pairs{1}{1}$ and the step adds $n+1$ to
both sides of the identity, which is what \cref{prop:4-1} asserts.
\end{proof}
\begin{equation}\label{eq:4-1}
  \int_{0}^{\infty}\frac{\vecn{x}^{s-1}}{e^{\vecn{x}}-1}\,d\vecn{x}
    = \Gamma(s)\zeta(s),\qquad \Re(s)>1 .
\end{equation}
\begin{align}
  \mathcal{L}[f](s) &= \int_{0}^{\infty} e^{-st} f(t)\,dt, \\
  \widehat{f}(\xi)  &= \int_{-\infty}^{\infty} f(x) e^{-2\pi i x\xi}\,dx .
\end{align}
\par\noindent\begin{center}
\begin{tabular}{lrrr}
\hline
Name & One & Two & Three \\
\hline
Alpha & 1 & 2 & 3 \\
Beta & 4 & 5 & 6 \\
Gamma & 7 & 8 & 9 \\
\hline
\end{tabular}
\end{center}
\begin{itemize}
  \item A first item with a formula $a^2+b^2=c^2$ in it.
  \item A second item that refers to \eqref{eq:4-1}.
  \item A third item, longer, so that the list breaks across a line and the
    paragraph builder has something to weigh when it chooses the break.
\end{itemize}
\subsection{A subsection}
Running text again, with \textit{italics}, \textbf{bold}, and
\texttt{typewriter}, and a formula $\alpha\beta\gamma\delta$ set inline so
that the math list machinery runs inside a paragraph. A displayed formula
follows:
\[ \lim_{n\to\infty}\left(1+\frac{1}{n}\right)^{n}=e . \]
And a little more prose to fill the page out, with an em-dash --- and
quotation marks ``like these'' and a footnote.\footnote{A footnote, which is
an insertion, so the page builder must split it.}
\subsection{A subsection}
Running text again, with \textit{italics}, \textbf{bold}, and
\texttt{typewriter}, and a formula $\alpha\beta\gamma\delta$ set inline so
that the math list machinery runs inside a paragraph. A displayed formula
follows:
\[ \lim_{n\to\infty}\left(1+\frac{1}{n}\right)^{n}=e . \]
And a little more prose to fill the page out, with an em-dash --- and
quotation marks ``like these'' and a footnote.\footnote{A footnote, which is
an insertion, so the page builder must split it.}
\subsection{A subsection}
Running text again, with \textit{italics}, \textbf{bold}, and
\texttt{typewriter}, and a formula $\alpha\beta\gamma\delta$ set inline so
that the math list machinery runs inside a paragraph. A displayed formula
follows:
\[ \lim_{n\to\infty}\left(1+\frac{1}{n}\right)^{n}=e . \]
And a little more prose to fill the page out, with an em-dash --- and
quotation marks ``like these'' and a footnote.\footnote{A footnote, which is
an insertion, so the page builder must split it.}
\section{Prose that fills lines and breaks pages}
\label{sec:prose4-2}
Ordinary running text is what a typesetting engine spends most of its effort
upon, so the training document supplies a great deal of it. The words are
chosen for their length and their hyphenation, with incomprehensibility and
antidisestablishmentarianism thrown in so that the patterns are exercised.
Cross references such as \cref{sec:prose4-2} and citations of a page such
as \pageref{sec:prose4-2} make the auxiliary machinery run, and a
paragraph long enough to break over several lines gives the line breaker
something to weigh.
\begin{proposition}\label{prop:4-2}
For every $n\in\mathbb{N}$ the sum $\sum_{k=1}^{n}k$ equals
$\tfrac{1}{2}n(n+1)$.
\end{proposition}
\begin{proof}
Induction on $n$; the base case is $\pairs{1}{1}$ and the step adds $n+1$ to
both sides of the identity, which is what \cref{prop:4-2} asserts.
\end{proof}
\begin{equation}\label{eq:4-2}
  \int_{0}^{\infty}\frac{\vecn{x}^{s-1}}{e^{\vecn{x}}-1}\,d\vecn{x}
    = \Gamma(s)\zeta(s),\qquad \Re(s)>1 .
\end{equation}
\begin{align}
  \mathcal{L}[f](s) &= \int_{0}^{\infty} e^{-st} f(t)\,dt, \\
  \widehat{f}(\xi)  &= \int_{-\infty}^{\infty} f(x) e^{-2\pi i x\xi}\,dx .
\end{align}
\par\noindent\begin{center}
\begin{tabular}{lrrr}
\hline
Name & One & Two & Three \\
\hline
Alpha & 1 & 2 & 3 \\
Beta & 4 & 5 & 6 \\
Gamma & 7 & 8 & 9 \\
\hline
\end{tabular}
\end{center}
\begin{itemize}
  \item A first item with a formula $a^2+b^2=c^2$ in it.
  \item A second item that refers to \eqref{eq:4-2}.
  \item A third item, longer, so that the list breaks across a line and the
    paragraph builder has something to weigh when it chooses the break.
\end{itemize}
\subsection{A subsection}
Running text again, with \textit{italics}, \textbf{bold}, and
\texttt{typewriter}, and a formula $\alpha\beta\gamma\delta$ set inline so
that the math list machinery runs inside a paragraph. A displayed formula
follows:
\[ \lim_{n\to\infty}\left(1+\frac{1}{n}\right)^{n}=e . \]
And a little more prose to fill the page out, with an em-dash --- and
quotation marks ``like these'' and a footnote.\footnote{A footnote, which is
an insertion, so the page builder must split it.}
\subsection{A subsection}
Running text again, with \textit{italics}, \textbf{bold}, and
\texttt{typewriter}, and a formula $\alpha\beta\gamma\delta$ set inline so
that the math list machinery runs inside a paragraph. A displayed formula
follows:
\[ \lim_{n\to\infty}\left(1+\frac{1}{n}\right)^{n}=e . \]
And a little more prose to fill the page out, with an em-dash --- and
quotation marks ``like these'' and a footnote.\footnote{A footnote, which is
an insertion, so the page builder must split it.}
\subsection{A subsection}
Running text again, with \textit{italics}, \textbf{bold}, and
\texttt{typewriter}, and a formula $\alpha\beta\gamma\delta$ set inline so
that the math list machinery runs inside a paragraph. A displayed formula
follows:
\[ \lim_{n\to\infty}\left(1+\frac{1}{n}\right)^{n}=e . \]
And a little more prose to fill the page out, with an em-dash --- and
quotation marks ``like these'' and a footnote.\footnote{A footnote, which is
an insertion, so the page builder must split it.}
\section{Prose that fills lines and breaks pages}
\label{sec:prose4-3}
Ordinary running text is what a typesetting engine spends most of its effort
upon, so the training document supplies a great deal of it. The words are
chosen for their length and their hyphenation, with incomprehensibility and
antidisestablishmentarianism thrown in so that the patterns are exercised.
Cross references such as \cref{sec:prose4-3} and citations of a page such
as \pageref{sec:prose4-3} make the auxiliary machinery run, and a
paragraph long enough to break over several lines gives the line breaker
something to weigh.
\begin{proposition}\label{prop:4-3}
For every $n\in\mathbb{N}$ the sum $\sum_{k=1}^{n}k$ equals
$\tfrac{1}{2}n(n+1)$.
\end{proposition}
\begin{proof}
Induction on $n$; the base case is $\pairs{1}{1}$ and the step adds $n+1$ to
both sides of the identity, which is what \cref{prop:4-3} asserts.
\end{proof}
\begin{equation}\label{eq:4-3}
  \int_{0}^{\infty}\frac{\vecn{x}^{s-1}}{e^{\vecn{x}}-1}\,d\vecn{x}
    = \Gamma(s)\zeta(s),\qquad \Re(s)>1 .
\end{equation}
\begin{align}
  \mathcal{L}[f](s) &= \int_{0}^{\infty} e^{-st} f(t)\,dt, \\
  \widehat{f}(\xi)  &= \int_{-\infty}^{\infty} f(x) e^{-2\pi i x\xi}\,dx .
\end{align}
\par\noindent\begin{center}
\begin{tabular}{lrrr}
\hline
Name & One & Two & Three \\
\hline
Alpha & 1 & 2 & 3 \\
Beta & 4 & 5 & 6 \\
Gamma & 7 & 8 & 9 \\
\hline
\end{tabular}
\end{center}
\begin{itemize}
  \item A first item with a formula $a^2+b^2=c^2$ in it.
  \item A second item that refers to \eqref{eq:4-3}.
  \item A third item, longer, so that the list breaks across a line and the
    paragraph builder has something to weigh when it chooses the break.
\end{itemize}
\subsection{A subsection}
Running text again, with \textit{italics}, \textbf{bold}, and
\texttt{typewriter}, and a formula $\alpha\beta\gamma\delta$ set inline so
that the math list machinery runs inside a paragraph. A displayed formula
follows:
\[ \lim_{n\to\infty}\left(1+\frac{1}{n}\right)^{n}=e . \]
And a little more prose to fill the page out, with an em-dash --- and
quotation marks ``like these'' and a footnote.\footnote{A footnote, which is
an insertion, so the page builder must split it.}
\subsection{A subsection}
Running text again, with \textit{italics}, \textbf{bold}, and
\texttt{typewriter}, and a formula $\alpha\beta\gamma\delta$ set inline so
that the math list machinery runs inside a paragraph. A displayed formula
follows:
\[ \lim_{n\to\infty}\left(1+\frac{1}{n}\right)^{n}=e . \]
And a little more prose to fill the page out, with an em-dash --- and
quotation marks ``like these'' and a footnote.\footnote{A footnote, which is
an insertion, so the page builder must split it.}
\subsection{A subsection}
Running text again, with \textit{italics}, \textbf{bold}, and
\texttt{typewriter}, and a formula $\alpha\beta\gamma\delta$ set inline so
that the math list machinery runs inside a paragraph. A displayed formula
follows:
\[ \lim_{n\to\infty}\left(1+\frac{1}{n}\right)^{n}=e . \]
And a little more prose to fill the page out, with an em-dash --- and
quotation marks ``like these'' and a footnote.\footnote{A footnote, which is
an insertion, so the page builder must split it.}
\section{Prose that fills lines and breaks pages}
\label{sec:prose4-4}
Ordinary running text is what a typesetting engine spends most of its effort
upon, so the training document supplies a great deal of it. The words are
chosen for their length and their hyphenation, with incomprehensibility and
antidisestablishmentarianism thrown in so that the patterns are exercised.
Cross references such as \cref{sec:prose4-4} and citations of a page such
as \pageref{sec:prose4-4} make the auxiliary machinery run, and a
paragraph long enough to break over several lines gives the line breaker
something to weigh.
\begin{proposition}\label{prop:4-4}
For every $n\in\mathbb{N}$ the sum $\sum_{k=1}^{n}k$ equals
$\tfrac{1}{2}n(n+1)$.
\end{proposition}
\begin{proof}
Induction on $n$; the base case is $\pairs{1}{1}$ and the step adds $n+1$ to
both sides of the identity, which is what \cref{prop:4-4} asserts.
\end{proof}
\begin{equation}\label{eq:4-4}
  \int_{0}^{\infty}\frac{\vecn{x}^{s-1}}{e^{\vecn{x}}-1}\,d\vecn{x}
    = \Gamma(s)\zeta(s),\qquad \Re(s)>1 .
\end{equation}
\begin{align}
  \mathcal{L}[f](s) &= \int_{0}^{\infty} e^{-st} f(t)\,dt, \\
  \widehat{f}(\xi)  &= \int_{-\infty}^{\infty} f(x) e^{-2\pi i x\xi}\,dx .
\end{align}
\par\noindent\begin{center}
\begin{tabular}{lrrr}
\hline
Name & One & Two & Three \\
\hline
Alpha & 1 & 2 & 3 \\
Beta & 4 & 5 & 6 \\
Gamma & 7 & 8 & 9 \\
\hline
\end{tabular}
\end{center}
\begin{itemize}
  \item A first item with a formula $a^2+b^2=c^2$ in it.
  \item A second item that refers to \eqref{eq:4-4}.
  \item A third item, longer, so that the list breaks across a line and the
    paragraph builder has something to weigh when it chooses the break.
\end{itemize}
\subsection{A subsection}
Running text again, with \textit{italics}, \textbf{bold}, and
\texttt{typewriter}, and a formula $\alpha\beta\gamma\delta$ set inline so
that the math list machinery runs inside a paragraph. A displayed formula
follows:
\[ \lim_{n\to\infty}\left(1+\frac{1}{n}\right)^{n}=e . \]
And a little more prose to fill the page out, with an em-dash --- and
quotation marks ``like these'' and a footnote.\footnote{A footnote, which is
an insertion, so the page builder must split it.}
\subsection{A subsection}
Running text again, with \textit{italics}, \textbf{bold}, and
\texttt{typewriter}, and a formula $\alpha\beta\gamma\delta$ set inline so
that the math list machinery runs inside a paragraph. A displayed formula
follows:
\[ \lim_{n\to\infty}\left(1+\frac{1}{n}\right)^{n}=e . \]
And a little more prose to fill the page out, with an em-dash --- and
quotation marks ``like these'' and a footnote.\footnote{A footnote, which is
an insertion, so the page builder must split it.}
\subsection{A subsection}
Running text again, with \textit{italics}, \textbf{bold}, and
\texttt{typewriter}, and a formula $\alpha\beta\gamma\delta$ set inline so
that the math list machinery runs inside a paragraph. A displayed formula
follows:
\[ \lim_{n\to\infty}\left(1+\frac{1}{n}\right)^{n}=e . \]
And a little more prose to fill the page out, with an em-dash --- and
quotation marks ``like these'' and a footnote.\footnote{A footnote, which is
an insertion, so the page builder must split it.}
\chapter{A chapter of ordinary matter, number 5}
\section{Prose that fills lines and breaks pages}
\label{sec:prose5-1}
Ordinary running text is what a typesetting engine spends most of its effort
upon, so the training document supplies a great deal of it. The words are
chosen for their length and their hyphenation, with incomprehensibility and
antidisestablishmentarianism thrown in so that the patterns are exercised.
Cross references such as \cref{sec:prose5-1} and citations of a page such
as \pageref{sec:prose5-1} make the auxiliary machinery run, and a
paragraph long enough to break over several lines gives the line breaker
something to weigh.
\begin{proposition}\label{prop:5-1}
For every $n\in\mathbb{N}$ the sum $\sum_{k=1}^{n}k$ equals
$\tfrac{1}{2}n(n+1)$.
\end{proposition}
\begin{proof}
Induction on $n$; the base case is $\pairs{1}{1}$ and the step adds $n+1$ to
both sides of the identity, which is what \cref{prop:5-1} asserts.
\end{proof}
\begin{equation}\label{eq:5-1}
  \int_{0}^{\infty}\frac{\vecn{x}^{s-1}}{e^{\vecn{x}}-1}\,d\vecn{x}
    = \Gamma(s)\zeta(s),\qquad \Re(s)>1 .
\end{equation}
\begin{align}
  \mathcal{L}[f](s) &= \int_{0}^{\infty} e^{-st} f(t)\,dt, \\
  \widehat{f}(\xi)  &= \int_{-\infty}^{\infty} f(x) e^{-2\pi i x\xi}\,dx .
\end{align}
\par\noindent\begin{center}
\begin{tabular}{lrrr}
\hline
Name & One & Two & Three \\
\hline
Alpha & 1 & 2 & 3 \\
Beta & 4 & 5 & 6 \\
Gamma & 7 & 8 & 9 \\
\hline
\end{tabular}
\end{center}
\begin{itemize}
  \item A first item with a formula $a^2+b^2=c^2$ in it.
  \item A second item that refers to \eqref{eq:5-1}.
  \item A third item, longer, so that the list breaks across a line and the
    paragraph builder has something to weigh when it chooses the break.
\end{itemize}
\subsection{A subsection}
Running text again, with \textit{italics}, \textbf{bold}, and
\texttt{typewriter}, and a formula $\alpha\beta\gamma\delta$ set inline so
that the math list machinery runs inside a paragraph. A displayed formula
follows:
\[ \lim_{n\to\infty}\left(1+\frac{1}{n}\right)^{n}=e . \]
And a little more prose to fill the page out, with an em-dash --- and
quotation marks ``like these'' and a footnote.\footnote{A footnote, which is
an insertion, so the page builder must split it.}
\subsection{A subsection}
Running text again, with \textit{italics}, \textbf{bold}, and
\texttt{typewriter}, and a formula $\alpha\beta\gamma\delta$ set inline so
that the math list machinery runs inside a paragraph. A displayed formula
follows:
\[ \lim_{n\to\infty}\left(1+\frac{1}{n}\right)^{n}=e . \]
And a little more prose to fill the page out, with an em-dash --- and
quotation marks ``like these'' and a footnote.\footnote{A footnote, which is
an insertion, so the page builder must split it.}
\subsection{A subsection}
Running text again, with \textit{italics}, \textbf{bold}, and
\texttt{typewriter}, and a formula $\alpha\beta\gamma\delta$ set inline so
that the math list machinery runs inside a paragraph. A displayed formula
follows:
\[ \lim_{n\to\infty}\left(1+\frac{1}{n}\right)^{n}=e . \]
And a little more prose to fill the page out, with an em-dash --- and
quotation marks ``like these'' and a footnote.\footnote{A footnote, which is
an insertion, so the page builder must split it.}
\section{Prose that fills lines and breaks pages}
\label{sec:prose5-2}
Ordinary running text is what a typesetting engine spends most of its effort
upon, so the training document supplies a great deal of it. The words are
chosen for their length and their hyphenation, with incomprehensibility and
antidisestablishmentarianism thrown in so that the patterns are exercised.
Cross references such as \cref{sec:prose5-2} and citations of a page such
as \pageref{sec:prose5-2} make the auxiliary machinery run, and a
paragraph long enough to break over several lines gives the line breaker
something to weigh.
\begin{proposition}\label{prop:5-2}
For every $n\in\mathbb{N}$ the sum $\sum_{k=1}^{n}k$ equals
$\tfrac{1}{2}n(n+1)$.
\end{proposition}
\begin{proof}
Induction on $n$; the base case is $\pairs{1}{1}$ and the step adds $n+1$ to
both sides of the identity, which is what \cref{prop:5-2} asserts.
\end{proof}
\begin{equation}\label{eq:5-2}
  \int_{0}^{\infty}\frac{\vecn{x}^{s-1}}{e^{\vecn{x}}-1}\,d\vecn{x}
    = \Gamma(s)\zeta(s),\qquad \Re(s)>1 .
\end{equation}
\begin{align}
  \mathcal{L}[f](s) &= \int_{0}^{\infty} e^{-st} f(t)\,dt, \\
  \widehat{f}(\xi)  &= \int_{-\infty}^{\infty} f(x) e^{-2\pi i x\xi}\,dx .
\end{align}
\par\noindent\begin{center}
\begin{tabular}{lrrr}
\hline
Name & One & Two & Three \\
\hline
Alpha & 1 & 2 & 3 \\
Beta & 4 & 5 & 6 \\
Gamma & 7 & 8 & 9 \\
\hline
\end{tabular}
\end{center}
\begin{itemize}
  \item A first item with a formula $a^2+b^2=c^2$ in it.
  \item A second item that refers to \eqref{eq:5-2}.
  \item A third item, longer, so that the list breaks across a line and the
    paragraph builder has something to weigh when it chooses the break.
\end{itemize}
\subsection{A subsection}
Running text again, with \textit{italics}, \textbf{bold}, and
\texttt{typewriter}, and a formula $\alpha\beta\gamma\delta$ set inline so
that the math list machinery runs inside a paragraph. A displayed formula
follows:
\[ \lim_{n\to\infty}\left(1+\frac{1}{n}\right)^{n}=e . \]
And a little more prose to fill the page out, with an em-dash --- and
quotation marks ``like these'' and a footnote.\footnote{A footnote, which is
an insertion, so the page builder must split it.}
\subsection{A subsection}
Running text again, with \textit{italics}, \textbf{bold}, and
\texttt{typewriter}, and a formula $\alpha\beta\gamma\delta$ set inline so
that the math list machinery runs inside a paragraph. A displayed formula
follows:
\[ \lim_{n\to\infty}\left(1+\frac{1}{n}\right)^{n}=e . \]
And a little more prose to fill the page out, with an em-dash --- and
quotation marks ``like these'' and a footnote.\footnote{A footnote, which is
an insertion, so the page builder must split it.}
\subsection{A subsection}
Running text again, with \textit{italics}, \textbf{bold}, and
\texttt{typewriter}, and a formula $\alpha\beta\gamma\delta$ set inline so
that the math list machinery runs inside a paragraph. A displayed formula
follows:
\[ \lim_{n\to\infty}\left(1+\frac{1}{n}\right)^{n}=e . \]
And a little more prose to fill the page out, with an em-dash --- and
quotation marks ``like these'' and a footnote.\footnote{A footnote, which is
an insertion, so the page builder must split it.}
\section{Prose that fills lines and breaks pages}
\label{sec:prose5-3}
Ordinary running text is what a typesetting engine spends most of its effort
upon, so the training document supplies a great deal of it. The words are
chosen for their length and their hyphenation, with incomprehensibility and
antidisestablishmentarianism thrown in so that the patterns are exercised.
Cross references such as \cref{sec:prose5-3} and citations of a page such
as \pageref{sec:prose5-3} make the auxiliary machinery run, and a
paragraph long enough to break over several lines gives the line breaker
something to weigh.
\begin{proposition}\label{prop:5-3}
For every $n\in\mathbb{N}$ the sum $\sum_{k=1}^{n}k$ equals
$\tfrac{1}{2}n(n+1)$.
\end{proposition}
\begin{proof}
Induction on $n$; the base case is $\pairs{1}{1}$ and the step adds $n+1$ to
both sides of the identity, which is what \cref{prop:5-3} asserts.
\end{proof}
\begin{equation}\label{eq:5-3}
  \int_{0}^{\infty}\frac{\vecn{x}^{s-1}}{e^{\vecn{x}}-1}\,d\vecn{x}
    = \Gamma(s)\zeta(s),\qquad \Re(s)>1 .
\end{equation}
\begin{align}
  \mathcal{L}[f](s) &= \int_{0}^{\infty} e^{-st} f(t)\,dt, \\
  \widehat{f}(\xi)  &= \int_{-\infty}^{\infty} f(x) e^{-2\pi i x\xi}\,dx .
\end{align}
\par\noindent\begin{center}
\begin{tabular}{lrrr}
\hline
Name & One & Two & Three \\
\hline
Alpha & 1 & 2 & 3 \\
Beta & 4 & 5 & 6 \\
Gamma & 7 & 8 & 9 \\
\hline
\end{tabular}
\end{center}
\begin{itemize}
  \item A first item with a formula $a^2+b^2=c^2$ in it.
  \item A second item that refers to \eqref{eq:5-3}.
  \item A third item, longer, so that the list breaks across a line and the
    paragraph builder has something to weigh when it chooses the break.
\end{itemize}
\subsection{A subsection}
Running text again, with \textit{italics}, \textbf{bold}, and
\texttt{typewriter}, and a formula $\alpha\beta\gamma\delta$ set inline so
that the math list machinery runs inside a paragraph. A displayed formula
follows:
\[ \lim_{n\to\infty}\left(1+\frac{1}{n}\right)^{n}=e . \]
And a little more prose to fill the page out, with an em-dash --- and
quotation marks ``like these'' and a footnote.\footnote{A footnote, which is
an insertion, so the page builder must split it.}
\subsection{A subsection}
Running text again, with \textit{italics}, \textbf{bold}, and
\texttt{typewriter}, and a formula $\alpha\beta\gamma\delta$ set inline so
that the math list machinery runs inside a paragraph. A displayed formula
follows:
\[ \lim_{n\to\infty}\left(1+\frac{1}{n}\right)^{n}=e . \]
And a little more prose to fill the page out, with an em-dash --- and
quotation marks ``like these'' and a footnote.\footnote{A footnote, which is
an insertion, so the page builder must split it.}
\subsection{A subsection}
Running text again, with \textit{italics}, \textbf{bold}, and
\texttt{typewriter}, and a formula $\alpha\beta\gamma\delta$ set inline so
that the math list machinery runs inside a paragraph. A displayed formula
follows:
\[ \lim_{n\to\infty}\left(1+\frac{1}{n}\right)^{n}=e . \]
And a little more prose to fill the page out, with an em-dash --- and
quotation marks ``like these'' and a footnote.\footnote{A footnote, which is
an insertion, so the page builder must split it.}
\section{Prose that fills lines and breaks pages}
\label{sec:prose5-4}
Ordinary running text is what a typesetting engine spends most of its effort
upon, so the training document supplies a great deal of it. The words are
chosen for their length and their hyphenation, with incomprehensibility and
antidisestablishmentarianism thrown in so that the patterns are exercised.
Cross references such as \cref{sec:prose5-4} and citations of a page such
as \pageref{sec:prose5-4} make the auxiliary machinery run, and a
paragraph long enough to break over several lines gives the line breaker
something to weigh.
\begin{proposition}\label{prop:5-4}
For every $n\in\mathbb{N}$ the sum $\sum_{k=1}^{n}k$ equals
$\tfrac{1}{2}n(n+1)$.
\end{proposition}
\begin{proof}
Induction on $n$; the base case is $\pairs{1}{1}$ and the step adds $n+1$ to
both sides of the identity, which is what \cref{prop:5-4} asserts.
\end{proof}
\begin{equation}\label{eq:5-4}
  \int_{0}^{\infty}\frac{\vecn{x}^{s-1}}{e^{\vecn{x}}-1}\,d\vecn{x}
    = \Gamma(s)\zeta(s),\qquad \Re(s)>1 .
\end{equation}
\begin{align}
  \mathcal{L}[f](s) &= \int_{0}^{\infty} e^{-st} f(t)\,dt, \\
  \widehat{f}(\xi)  &= \int_{-\infty}^{\infty} f(x) e^{-2\pi i x\xi}\,dx .
\end{align}
\par\noindent\begin{center}
\begin{tabular}{lrrr}
\hline
Name & One & Two & Three \\
\hline
Alpha & 1 & 2 & 3 \\
Beta & 4 & 5 & 6 \\
Gamma & 7 & 8 & 9 \\
\hline
\end{tabular}
\end{center}
\begin{itemize}
  \item A first item with a formula $a^2+b^2=c^2$ in it.
  \item A second item that refers to \eqref{eq:5-4}.
  \item A third item, longer, so that the list breaks across a line and the
    paragraph builder has something to weigh when it chooses the break.
\end{itemize}
\subsection{A subsection}
Running text again, with \textit{italics}, \textbf{bold}, and
\texttt{typewriter}, and a formula $\alpha\beta\gamma\delta$ set inline so
that the math list machinery runs inside a paragraph. A displayed formula
follows:
\[ \lim_{n\to\infty}\left(1+\frac{1}{n}\right)^{n}=e . \]
And a little more prose to fill the page out, with an em-dash --- and
quotation marks ``like these'' and a footnote.\footnote{A footnote, which is
an insertion, so the page builder must split it.}
\subsection{A subsection}
Running text again, with \textit{italics}, \textbf{bold}, and
\texttt{typewriter}, and a formula $\alpha\beta\gamma\delta$ set inline so
that the math list machinery runs inside a paragraph. A displayed formula
follows:
\[ \lim_{n\to\infty}\left(1+\frac{1}{n}\right)^{n}=e . \]
And a little more prose to fill the page out, with an em-dash --- and
quotation marks ``like these'' and a footnote.\footnote{A footnote, which is
an insertion, so the page builder must split it.}
\subsection{A subsection}
Running text again, with \textit{italics}, \textbf{bold}, and
\texttt{typewriter}, and a formula $\alpha\beta\gamma\delta$ set inline so
that the math list machinery runs inside a paragraph. A displayed formula
follows:
\[ \lim_{n\to\infty}\left(1+\frac{1}{n}\right)^{n}=e . \]
And a little more prose to fill the page out, with an em-dash --- and
quotation marks ``like these'' and a footnote.\footnote{A footnote, which is
an insertion, so the page builder must split it.}
\chapter{A chapter of ordinary matter, number 6}
\section{Prose that fills lines and breaks pages}
\label{sec:prose6-1}
Ordinary running text is what a typesetting engine spends most of its effort
upon, so the training document supplies a great deal of it. The words are
chosen for their length and their hyphenation, with incomprehensibility and
antidisestablishmentarianism thrown in so that the patterns are exercised.
Cross references such as \cref{sec:prose6-1} and citations of a page such
as \pageref{sec:prose6-1} make the auxiliary machinery run, and a
paragraph long enough to break over several lines gives the line breaker
something to weigh.
\begin{proposition}\label{prop:6-1}
For every $n\in\mathbb{N}$ the sum $\sum_{k=1}^{n}k$ equals
$\tfrac{1}{2}n(n+1)$.
\end{proposition}
\begin{proof}
Induction on $n$; the base case is $\pairs{1}{1}$ and the step adds $n+1$ to
both sides of the identity, which is what \cref{prop:6-1} asserts.
\end{proof}
\begin{equation}\label{eq:6-1}
  \int_{0}^{\infty}\frac{\vecn{x}^{s-1}}{e^{\vecn{x}}-1}\,d\vecn{x}
    = \Gamma(s)\zeta(s),\qquad \Re(s)>1 .
\end{equation}
\begin{align}
  \mathcal{L}[f](s) &= \int_{0}^{\infty} e^{-st} f(t)\,dt, \\
  \widehat{f}(\xi)  &= \int_{-\infty}^{\infty} f(x) e^{-2\pi i x\xi}\,dx .
\end{align}
\par\noindent\begin{center}
\begin{tabular}{lrrr}
\hline
Name & One & Two & Three \\
\hline
Alpha & 1 & 2 & 3 \\
Beta & 4 & 5 & 6 \\
Gamma & 7 & 8 & 9 \\
\hline
\end{tabular}
\end{center}
\begin{itemize}
  \item A first item with a formula $a^2+b^2=c^2$ in it.
  \item A second item that refers to \eqref{eq:6-1}.
  \item A third item, longer, so that the list breaks across a line and the
    paragraph builder has something to weigh when it chooses the break.
\end{itemize}
\subsection{A subsection}
Running text again, with \textit{italics}, \textbf{bold}, and
\texttt{typewriter}, and a formula $\alpha\beta\gamma\delta$ set inline so
that the math list machinery runs inside a paragraph. A displayed formula
follows:
\[ \lim_{n\to\infty}\left(1+\frac{1}{n}\right)^{n}=e . \]
And a little more prose to fill the page out, with an em-dash --- and
quotation marks ``like these'' and a footnote.\footnote{A footnote, which is
an insertion, so the page builder must split it.}
\subsection{A subsection}
Running text again, with \textit{italics}, \textbf{bold}, and
\texttt{typewriter}, and a formula $\alpha\beta\gamma\delta$ set inline so
that the math list machinery runs inside a paragraph. A displayed formula
follows:
\[ \lim_{n\to\infty}\left(1+\frac{1}{n}\right)^{n}=e . \]
And a little more prose to fill the page out, with an em-dash --- and
quotation marks ``like these'' and a footnote.\footnote{A footnote, which is
an insertion, so the page builder must split it.}
\subsection{A subsection}
Running text again, with \textit{italics}, \textbf{bold}, and
\texttt{typewriter}, and a formula $\alpha\beta\gamma\delta$ set inline so
that the math list machinery runs inside a paragraph. A displayed formula
follows:
\[ \lim_{n\to\infty}\left(1+\frac{1}{n}\right)^{n}=e . \]
And a little more prose to fill the page out, with an em-dash --- and
quotation marks ``like these'' and a footnote.\footnote{A footnote, which is
an insertion, so the page builder must split it.}
\section{Prose that fills lines and breaks pages}
\label{sec:prose6-2}
Ordinary running text is what a typesetting engine spends most of its effort
upon, so the training document supplies a great deal of it. The words are
chosen for their length and their hyphenation, with incomprehensibility and
antidisestablishmentarianism thrown in so that the patterns are exercised.
Cross references such as \cref{sec:prose6-2} and citations of a page such
as \pageref{sec:prose6-2} make the auxiliary machinery run, and a
paragraph long enough to break over several lines gives the line breaker
something to weigh.
\begin{proposition}\label{prop:6-2}
For every $n\in\mathbb{N}$ the sum $\sum_{k=1}^{n}k$ equals
$\tfrac{1}{2}n(n+1)$.
\end{proposition}
\begin{proof}
Induction on $n$; the base case is $\pairs{1}{1}$ and the step adds $n+1$ to
both sides of the identity, which is what \cref{prop:6-2} asserts.
\end{proof}
\begin{equation}\label{eq:6-2}
  \int_{0}^{\infty}\frac{\vecn{x}^{s-1}}{e^{\vecn{x}}-1}\,d\vecn{x}
    = \Gamma(s)\zeta(s),\qquad \Re(s)>1 .
\end{equation}
\begin{align}
  \mathcal{L}[f](s) &= \int_{0}^{\infty} e^{-st} f(t)\,dt, \\
  \widehat{f}(\xi)  &= \int_{-\infty}^{\infty} f(x) e^{-2\pi i x\xi}\,dx .
\end{align}
\par\noindent\begin{center}
\begin{tabular}{lrrr}
\hline
Name & One & Two & Three \\
\hline
Alpha & 1 & 2 & 3 \\
Beta & 4 & 5 & 6 \\
Gamma & 7 & 8 & 9 \\
\hline
\end{tabular}
\end{center}
\begin{itemize}
  \item A first item with a formula $a^2+b^2=c^2$ in it.
  \item A second item that refers to \eqref{eq:6-2}.
  \item A third item, longer, so that the list breaks across a line and the
    paragraph builder has something to weigh when it chooses the break.
\end{itemize}
\subsection{A subsection}
Running text again, with \textit{italics}, \textbf{bold}, and
\texttt{typewriter}, and a formula $\alpha\beta\gamma\delta$ set inline so
that the math list machinery runs inside a paragraph. A displayed formula
follows:
\[ \lim_{n\to\infty}\left(1+\frac{1}{n}\right)^{n}=e . \]
And a little more prose to fill the page out, with an em-dash --- and
quotation marks ``like these'' and a footnote.\footnote{A footnote, which is
an insertion, so the page builder must split it.}
\subsection{A subsection}
Running text again, with \textit{italics}, \textbf{bold}, and
\texttt{typewriter}, and a formula $\alpha\beta\gamma\delta$ set inline so
that the math list machinery runs inside a paragraph. A displayed formula
follows:
\[ \lim_{n\to\infty}\left(1+\frac{1}{n}\right)^{n}=e . \]
And a little more prose to fill the page out, with an em-dash --- and
quotation marks ``like these'' and a footnote.\footnote{A footnote, which is
an insertion, so the page builder must split it.}
\subsection{A subsection}
Running text again, with \textit{italics}, \textbf{bold}, and
\texttt{typewriter}, and a formula $\alpha\beta\gamma\delta$ set inline so
that the math list machinery runs inside a paragraph. A displayed formula
follows:
\[ \lim_{n\to\infty}\left(1+\frac{1}{n}\right)^{n}=e . \]
And a little more prose to fill the page out, with an em-dash --- and
quotation marks ``like these'' and a footnote.\footnote{A footnote, which is
an insertion, so the page builder must split it.}
\section{Prose that fills lines and breaks pages}
\label{sec:prose6-3}
Ordinary running text is what a typesetting engine spends most of its effort
upon, so the training document supplies a great deal of it. The words are
chosen for their length and their hyphenation, with incomprehensibility and
antidisestablishmentarianism thrown in so that the patterns are exercised.
Cross references such as \cref{sec:prose6-3} and citations of a page such
as \pageref{sec:prose6-3} make the auxiliary machinery run, and a
paragraph long enough to break over several lines gives the line breaker
something to weigh.
\begin{proposition}\label{prop:6-3}
For every $n\in\mathbb{N}$ the sum $\sum_{k=1}^{n}k$ equals
$\tfrac{1}{2}n(n+1)$.
\end{proposition}
\begin{proof}
Induction on $n$; the base case is $\pairs{1}{1}$ and the step adds $n+1$ to
both sides of the identity, which is what \cref{prop:6-3} asserts.
\end{proof}
\begin{equation}\label{eq:6-3}
  \int_{0}^{\infty}\frac{\vecn{x}^{s-1}}{e^{\vecn{x}}-1}\,d\vecn{x}
    = \Gamma(s)\zeta(s),\qquad \Re(s)>1 .
\end{equation}
\begin{align}
  \mathcal{L}[f](s) &= \int_{0}^{\infty} e^{-st} f(t)\,dt, \\
  \widehat{f}(\xi)  &= \int_{-\infty}^{\infty} f(x) e^{-2\pi i x\xi}\,dx .
\end{align}
\par\noindent\begin{center}
\begin{tabular}{lrrr}
\hline
Name & One & Two & Three \\
\hline
Alpha & 1 & 2 & 3 \\
Beta & 4 & 5 & 6 \\
Gamma & 7 & 8 & 9 \\
\hline
\end{tabular}
\end{center}
\begin{itemize}
  \item A first item with a formula $a^2+b^2=c^2$ in it.
  \item A second item that refers to \eqref{eq:6-3}.
  \item A third item, longer, so that the list breaks across a line and the
    paragraph builder has something to weigh when it chooses the break.
\end{itemize}
\subsection{A subsection}
Running text again, with \textit{italics}, \textbf{bold}, and
\texttt{typewriter}, and a formula $\alpha\beta\gamma\delta$ set inline so
that the math list machinery runs inside a paragraph. A displayed formula
follows:
\[ \lim_{n\to\infty}\left(1+\frac{1}{n}\right)^{n}=e . \]
And a little more prose to fill the page out, with an em-dash --- and
quotation marks ``like these'' and a footnote.\footnote{A footnote, which is
an insertion, so the page builder must split it.}
\subsection{A subsection}
Running text again, with \textit{italics}, \textbf{bold}, and
\texttt{typewriter}, and a formula $\alpha\beta\gamma\delta$ set inline so
that the math list machinery runs inside a paragraph. A displayed formula
follows:
\[ \lim_{n\to\infty}\left(1+\frac{1}{n}\right)^{n}=e . \]
And a little more prose to fill the page out, with an em-dash --- and
quotation marks ``like these'' and a footnote.\footnote{A footnote, which is
an insertion, so the page builder must split it.}
\subsection{A subsection}
Running text again, with \textit{italics}, \textbf{bold}, and
\texttt{typewriter}, and a formula $\alpha\beta\gamma\delta$ set inline so
that the math list machinery runs inside a paragraph. A displayed formula
follows:
\[ \lim_{n\to\infty}\left(1+\frac{1}{n}\right)^{n}=e . \]
And a little more prose to fill the page out, with an em-dash --- and
quotation marks ``like these'' and a footnote.\footnote{A footnote, which is
an insertion, so the page builder must split it.}
\section{Prose that fills lines and breaks pages}
\label{sec:prose6-4}
Ordinary running text is what a typesetting engine spends most of its effort
upon, so the training document supplies a great deal of it. The words are
chosen for their length and their hyphenation, with incomprehensibility and
antidisestablishmentarianism thrown in so that the patterns are exercised.
Cross references such as \cref{sec:prose6-4} and citations of a page such
as \pageref{sec:prose6-4} make the auxiliary machinery run, and a
paragraph long enough to break over several lines gives the line breaker
something to weigh.
\begin{proposition}\label{prop:6-4}
For every $n\in\mathbb{N}$ the sum $\sum_{k=1}^{n}k$ equals
$\tfrac{1}{2}n(n+1)$.
\end{proposition}
\begin{proof}
Induction on $n$; the base case is $\pairs{1}{1}$ and the step adds $n+1$ to
both sides of the identity, which is what \cref{prop:6-4} asserts.
\end{proof}
\begin{equation}\label{eq:6-4}
  \int_{0}^{\infty}\frac{\vecn{x}^{s-1}}{e^{\vecn{x}}-1}\,d\vecn{x}
    = \Gamma(s)\zeta(s),\qquad \Re(s)>1 .
\end{equation}
\begin{align}
  \mathcal{L}[f](s) &= \int_{0}^{\infty} e^{-st} f(t)\,dt, \\
  \widehat{f}(\xi)  &= \int_{-\infty}^{\infty} f(x) e^{-2\pi i x\xi}\,dx .
\end{align}
\par\noindent\begin{center}
\begin{tabular}{lrrr}
\hline
Name & One & Two & Three \\
\hline
Alpha & 1 & 2 & 3 \\
Beta & 4 & 5 & 6 \\
Gamma & 7 & 8 & 9 \\
\hline
\end{tabular}
\end{center}
\begin{itemize}
  \item A first item with a formula $a^2+b^2=c^2$ in it.
  \item A second item that refers to \eqref{eq:6-4}.
  \item A third item, longer, so that the list breaks across a line and the
    paragraph builder has something to weigh when it chooses the break.
\end{itemize}
\subsection{A subsection}
Running text again, with \textit{italics}, \textbf{bold}, and
\texttt{typewriter}, and a formula $\alpha\beta\gamma\delta$ set inline so
that the math list machinery runs inside a paragraph. A displayed formula
follows:
\[ \lim_{n\to\infty}\left(1+\frac{1}{n}\right)^{n}=e . \]
And a little more prose to fill the page out, with an em-dash --- and
quotation marks ``like these'' and a footnote.\footnote{A footnote, which is
an insertion, so the page builder must split it.}
\subsection{A subsection}
Running text again, with \textit{italics}, \textbf{bold}, and
\texttt{typewriter}, and a formula $\alpha\beta\gamma\delta$ set inline so
that the math list machinery runs inside a paragraph. A displayed formula
follows:
\[ \lim_{n\to\infty}\left(1+\frac{1}{n}\right)^{n}=e . \]
And a little more prose to fill the page out, with an em-dash --- and
quotation marks ``like these'' and a footnote.\footnote{A footnote, which is
an insertion, so the page builder must split it.}
\subsection{A subsection}
Running text again, with \textit{italics}, \textbf{bold}, and
\texttt{typewriter}, and a formula $\alpha\beta\gamma\delta$ set inline so
that the math list machinery runs inside a paragraph. A displayed formula
follows:
\[ \lim_{n\to\infty}\left(1+\frac{1}{n}\right)^{n}=e . \]
And a little more prose to fill the page out, with an em-dash --- and
quotation marks ``like these'' and a footnote.\footnote{A footnote, which is
an insertion, so the page builder must split it.}
\chapter{A chapter of ordinary matter, number 7}
\section{Prose that fills lines and breaks pages}
\label{sec:prose7-1}
Ordinary running text is what a typesetting engine spends most of its effort
upon, so the training document supplies a great deal of it. The words are
chosen for their length and their hyphenation, with incomprehensibility and
antidisestablishmentarianism thrown in so that the patterns are exercised.
Cross references such as \cref{sec:prose7-1} and citations of a page such
as \pageref{sec:prose7-1} make the auxiliary machinery run, and a
paragraph long enough to break over several lines gives the line breaker
something to weigh.
\begin{proposition}\label{prop:7-1}
For every $n\in\mathbb{N}$ the sum $\sum_{k=1}^{n}k$ equals
$\tfrac{1}{2}n(n+1)$.
\end{proposition}
\begin{proof}
Induction on $n$; the base case is $\pairs{1}{1}$ and the step adds $n+1$ to
both sides of the identity, which is what \cref{prop:7-1} asserts.
\end{proof}
\begin{equation}\label{eq:7-1}
  \int_{0}^{\infty}\frac{\vecn{x}^{s-1}}{e^{\vecn{x}}-1}\,d\vecn{x}
    = \Gamma(s)\zeta(s),\qquad \Re(s)>1 .
\end{equation}
\begin{align}
  \mathcal{L}[f](s) &= \int_{0}^{\infty} e^{-st} f(t)\,dt, \\
  \widehat{f}(\xi)  &= \int_{-\infty}^{\infty} f(x) e^{-2\pi i x\xi}\,dx .
\end{align}
\par\noindent\begin{center}
\begin{tabular}{lrrr}
\hline
Name & One & Two & Three \\
\hline
Alpha & 1 & 2 & 3 \\
Beta & 4 & 5 & 6 \\
Gamma & 7 & 8 & 9 \\
\hline
\end{tabular}
\end{center}
\begin{itemize}
  \item A first item with a formula $a^2+b^2=c^2$ in it.
  \item A second item that refers to \eqref{eq:7-1}.
  \item A third item, longer, so that the list breaks across a line and the
    paragraph builder has something to weigh when it chooses the break.
\end{itemize}
\subsection{A subsection}
Running text again, with \textit{italics}, \textbf{bold}, and
\texttt{typewriter}, and a formula $\alpha\beta\gamma\delta$ set inline so
that the math list machinery runs inside a paragraph. A displayed formula
follows:
\[ \lim_{n\to\infty}\left(1+\frac{1}{n}\right)^{n}=e . \]
And a little more prose to fill the page out, with an em-dash --- and
quotation marks ``like these'' and a footnote.\footnote{A footnote, which is
an insertion, so the page builder must split it.}
\subsection{A subsection}
Running text again, with \textit{italics}, \textbf{bold}, and
\texttt{typewriter}, and a formula $\alpha\beta\gamma\delta$ set inline so
that the math list machinery runs inside a paragraph. A displayed formula
follows:
\[ \lim_{n\to\infty}\left(1+\frac{1}{n}\right)^{n}=e . \]
And a little more prose to fill the page out, with an em-dash --- and
quotation marks ``like these'' and a footnote.\footnote{A footnote, which is
an insertion, so the page builder must split it.}
\subsection{A subsection}
Running text again, with \textit{italics}, \textbf{bold}, and
\texttt{typewriter}, and a formula $\alpha\beta\gamma\delta$ set inline so
that the math list machinery runs inside a paragraph. A displayed formula
follows:
\[ \lim_{n\to\infty}\left(1+\frac{1}{n}\right)^{n}=e . \]
And a little more prose to fill the page out, with an em-dash --- and
quotation marks ``like these'' and a footnote.\footnote{A footnote, which is
an insertion, so the page builder must split it.}
\section{Prose that fills lines and breaks pages}
\label{sec:prose7-2}
Ordinary running text is what a typesetting engine spends most of its effort
upon, so the training document supplies a great deal of it. The words are
chosen for their length and their hyphenation, with incomprehensibility and
antidisestablishmentarianism thrown in so that the patterns are exercised.
Cross references such as \cref{sec:prose7-2} and citations of a page such
as \pageref{sec:prose7-2} make the auxiliary machinery run, and a
paragraph long enough to break over several lines gives the line breaker
something to weigh.
\begin{proposition}\label{prop:7-2}
For every $n\in\mathbb{N}$ the sum $\sum_{k=1}^{n}k$ equals
$\tfrac{1}{2}n(n+1)$.
\end{proposition}
\begin{proof}
Induction on $n$; the base case is $\pairs{1}{1}$ and the step adds $n+1$ to
both sides of the identity, which is what \cref{prop:7-2} asserts.
\end{proof}
\begin{equation}\label{eq:7-2}
  \int_{0}^{\infty}\frac{\vecn{x}^{s-1}}{e^{\vecn{x}}-1}\,d\vecn{x}
    = \Gamma(s)\zeta(s),\qquad \Re(s)>1 .
\end{equation}
\begin{align}
  \mathcal{L}[f](s) &= \int_{0}^{\infty} e^{-st} f(t)\,dt, \\
  \widehat{f}(\xi)  &= \int_{-\infty}^{\infty} f(x) e^{-2\pi i x\xi}\,dx .
\end{align}
\par\noindent\begin{center}
\begin{tabular}{lrrr}
\hline
Name & One & Two & Three \\
\hline
Alpha & 1 & 2 & 3 \\
Beta & 4 & 5 & 6 \\
Gamma & 7 & 8 & 9 \\
\hline
\end{tabular}
\end{center}
\begin{itemize}
  \item A first item with a formula $a^2+b^2=c^2$ in it.
  \item A second item that refers to \eqref{eq:7-2}.
  \item A third item, longer, so that the list breaks across a line and the
    paragraph builder has something to weigh when it chooses the break.
\end{itemize}
\subsection{A subsection}
Running text again, with \textit{italics}, \textbf{bold}, and
\texttt{typewriter}, and a formula $\alpha\beta\gamma\delta$ set inline so
that the math list machinery runs inside a paragraph. A displayed formula
follows:
\[ \lim_{n\to\infty}\left(1+\frac{1}{n}\right)^{n}=e . \]
And a little more prose to fill the page out, with an em-dash --- and
quotation marks ``like these'' and a footnote.\footnote{A footnote, which is
an insertion, so the page builder must split it.}
\subsection{A subsection}
Running text again, with \textit{italics}, \textbf{bold}, and
\texttt{typewriter}, and a formula $\alpha\beta\gamma\delta$ set inline so
that the math list machinery runs inside a paragraph. A displayed formula
follows:
\[ \lim_{n\to\infty}\left(1+\frac{1}{n}\right)^{n}=e . \]
And a little more prose to fill the page out, with an em-dash --- and
quotation marks ``like these'' and a footnote.\footnote{A footnote, which is
an insertion, so the page builder must split it.}
\subsection{A subsection}
Running text again, with \textit{italics}, \textbf{bold}, and
\texttt{typewriter}, and a formula $\alpha\beta\gamma\delta$ set inline so
that the math list machinery runs inside a paragraph. A displayed formula
follows:
\[ \lim_{n\to\infty}\left(1+\frac{1}{n}\right)^{n}=e . \]
And a little more prose to fill the page out, with an em-dash --- and
quotation marks ``like these'' and a footnote.\footnote{A footnote, which is
an insertion, so the page builder must split it.}
\section{Prose that fills lines and breaks pages}
\label{sec:prose7-3}
Ordinary running text is what a typesetting engine spends most of its effort
upon, so the training document supplies a great deal of it. The words are
chosen for their length and their hyphenation, with incomprehensibility and
antidisestablishmentarianism thrown in so that the patterns are exercised.
Cross references such as \cref{sec:prose7-3} and citations of a page such
as \pageref{sec:prose7-3} make the auxiliary machinery run, and a
paragraph long enough to break over several lines gives the line breaker
something to weigh.
\begin{proposition}\label{prop:7-3}
For every $n\in\mathbb{N}$ the sum $\sum_{k=1}^{n}k$ equals
$\tfrac{1}{2}n(n+1)$.
\end{proposition}
\begin{proof}
Induction on $n$; the base case is $\pairs{1}{1}$ and the step adds $n+1$ to
both sides of the identity, which is what \cref{prop:7-3} asserts.
\end{proof}
\begin{equation}\label{eq:7-3}
  \int_{0}^{\infty}\frac{\vecn{x}^{s-1}}{e^{\vecn{x}}-1}\,d\vecn{x}
    = \Gamma(s)\zeta(s),\qquad \Re(s)>1 .
\end{equation}
\begin{align}
  \mathcal{L}[f](s) &= \int_{0}^{\infty} e^{-st} f(t)\,dt, \\
  \widehat{f}(\xi)  &= \int_{-\infty}^{\infty} f(x) e^{-2\pi i x\xi}\,dx .
\end{align}
\par\noindent\begin{center}
\begin{tabular}{lrrr}
\hline
Name & One & Two & Three \\
\hline
Alpha & 1 & 2 & 3 \\
Beta & 4 & 5 & 6 \\
Gamma & 7 & 8 & 9 \\
\hline
\end{tabular}
\end{center}
\begin{itemize}
  \item A first item with a formula $a^2+b^2=c^2$ in it.
  \item A second item that refers to \eqref{eq:7-3}.
  \item A third item, longer, so that the list breaks across a line and the
    paragraph builder has something to weigh when it chooses the break.
\end{itemize}
\subsection{A subsection}
Running text again, with \textit{italics}, \textbf{bold}, and
\texttt{typewriter}, and a formula $\alpha\beta\gamma\delta$ set inline so
that the math list machinery runs inside a paragraph. A displayed formula
follows:
\[ \lim_{n\to\infty}\left(1+\frac{1}{n}\right)^{n}=e . \]
And a little more prose to fill the page out, with an em-dash --- and
quotation marks ``like these'' and a footnote.\footnote{A footnote, which is
an insertion, so the page builder must split it.}
\subsection{A subsection}
Running text again, with \textit{italics}, \textbf{bold}, and
\texttt{typewriter}, and a formula $\alpha\beta\gamma\delta$ set inline so
that the math list machinery runs inside a paragraph. A displayed formula
follows:
\[ \lim_{n\to\infty}\left(1+\frac{1}{n}\right)^{n}=e . \]
And a little more prose to fill the page out, with an em-dash --- and
quotation marks ``like these'' and a footnote.\footnote{A footnote, which is
an insertion, so the page builder must split it.}
\subsection{A subsection}
Running text again, with \textit{italics}, \textbf{bold}, and
\texttt{typewriter}, and a formula $\alpha\beta\gamma\delta$ set inline so
that the math list machinery runs inside a paragraph. A displayed formula
follows:
\[ \lim_{n\to\infty}\left(1+\frac{1}{n}\right)^{n}=e . \]
And a little more prose to fill the page out, with an em-dash --- and
quotation marks ``like these'' and a footnote.\footnote{A footnote, which is
an insertion, so the page builder must split it.}
\section{Prose that fills lines and breaks pages}
\label{sec:prose7-4}
Ordinary running text is what a typesetting engine spends most of its effort
upon, so the training document supplies a great deal of it. The words are
chosen for their length and their hyphenation, with incomprehensibility and
antidisestablishmentarianism thrown in so that the patterns are exercised.
Cross references such as \cref{sec:prose7-4} and citations of a page such
as \pageref{sec:prose7-4} make the auxiliary machinery run, and a
paragraph long enough to break over several lines gives the line breaker
something to weigh.
\begin{proposition}\label{prop:7-4}
For every $n\in\mathbb{N}$ the sum $\sum_{k=1}^{n}k$ equals
$\tfrac{1}{2}n(n+1)$.
\end{proposition}
\begin{proof}
Induction on $n$; the base case is $\pairs{1}{1}$ and the step adds $n+1$ to
both sides of the identity, which is what \cref{prop:7-4} asserts.
\end{proof}
\begin{equation}\label{eq:7-4}
  \int_{0}^{\infty}\frac{\vecn{x}^{s-1}}{e^{\vecn{x}}-1}\,d\vecn{x}
    = \Gamma(s)\zeta(s),\qquad \Re(s)>1 .
\end{equation}
\begin{align}
  \mathcal{L}[f](s) &= \int_{0}^{\infty} e^{-st} f(t)\,dt, \\
  \widehat{f}(\xi)  &= \int_{-\infty}^{\infty} f(x) e^{-2\pi i x\xi}\,dx .
\end{align}
\par\noindent\begin{center}
\begin{tabular}{lrrr}
\hline
Name & One & Two & Three \\
\hline
Alpha & 1 & 2 & 3 \\
Beta & 4 & 5 & 6 \\
Gamma & 7 & 8 & 9 \\
\hline
\end{tabular}
\end{center}
\begin{itemize}
  \item A first item with a formula $a^2+b^2=c^2$ in it.
  \item A second item that refers to \eqref{eq:7-4}.
  \item A third item, longer, so that the list breaks across a line and the
    paragraph builder has something to weigh when it chooses the break.
\end{itemize}
\subsection{A subsection}
Running text again, with \textit{italics}, \textbf{bold}, and
\texttt{typewriter}, and a formula $\alpha\beta\gamma\delta$ set inline so
that the math list machinery runs inside a paragraph. A displayed formula
follows:
\[ \lim_{n\to\infty}\left(1+\frac{1}{n}\right)^{n}=e . \]
And a little more prose to fill the page out, with an em-dash --- and
quotation marks ``like these'' and a footnote.\footnote{A footnote, which is
an insertion, so the page builder must split it.}
\subsection{A subsection}
Running text again, with \textit{italics}, \textbf{bold}, and
\texttt{typewriter}, and a formula $\alpha\beta\gamma\delta$ set inline so
that the math list machinery runs inside a paragraph. A displayed formula
follows:
\[ \lim_{n\to\infty}\left(1+\frac{1}{n}\right)^{n}=e . \]
And a little more prose to fill the page out, with an em-dash --- and
quotation marks ``like these'' and a footnote.\footnote{A footnote, which is
an insertion, so the page builder must split it.}
\subsection{A subsection}
Running text again, with \textit{italics}, \textbf{bold}, and
\texttt{typewriter}, and a formula $\alpha\beta\gamma\delta$ set inline so
that the math list machinery runs inside a paragraph. A displayed formula
follows:
\[ \lim_{n\to\infty}\left(1+\frac{1}{n}\right)^{n}=e . \]
And a little more prose to fill the page out, with an em-dash --- and
quotation marks ``like these'' and a footnote.\footnote{A footnote, which is
an insertion, so the page builder must split it.}
\chapter{A chapter of ordinary matter, number 8}
\section{Prose that fills lines and breaks pages}
\label{sec:prose8-1}
Ordinary running text is what a typesetting engine spends most of its effort
upon, so the training document supplies a great deal of it. The words are
chosen for their length and their hyphenation, with incomprehensibility and
antidisestablishmentarianism thrown in so that the patterns are exercised.
Cross references such as \cref{sec:prose8-1} and citations of a page such
as \pageref{sec:prose8-1} make the auxiliary machinery run, and a
paragraph long enough to break over several lines gives the line breaker
something to weigh.
\begin{proposition}\label{prop:8-1}
For every $n\in\mathbb{N}$ the sum $\sum_{k=1}^{n}k$ equals
$\tfrac{1}{2}n(n+1)$.
\end{proposition}
\begin{proof}
Induction on $n$; the base case is $\pairs{1}{1}$ and the step adds $n+1$ to
both sides of the identity, which is what \cref{prop:8-1} asserts.
\end{proof}
\begin{equation}\label{eq:8-1}
  \int_{0}^{\infty}\frac{\vecn{x}^{s-1}}{e^{\vecn{x}}-1}\,d\vecn{x}
    = \Gamma(s)\zeta(s),\qquad \Re(s)>1 .
\end{equation}
\begin{align}
  \mathcal{L}[f](s) &= \int_{0}^{\infty} e^{-st} f(t)\,dt, \\
  \widehat{f}(\xi)  &= \int_{-\infty}^{\infty} f(x) e^{-2\pi i x\xi}\,dx .
\end{align}
\par\noindent\begin{center}
\begin{tabular}{lrrr}
\hline
Name & One & Two & Three \\
\hline
Alpha & 1 & 2 & 3 \\
Beta & 4 & 5 & 6 \\
Gamma & 7 & 8 & 9 \\
\hline
\end{tabular}
\end{center}
\begin{itemize}
  \item A first item with a formula $a^2+b^2=c^2$ in it.
  \item A second item that refers to \eqref{eq:8-1}.
  \item A third item, longer, so that the list breaks across a line and the
    paragraph builder has something to weigh when it chooses the break.
\end{itemize}
\subsection{A subsection}
Running text again, with \textit{italics}, \textbf{bold}, and
\texttt{typewriter}, and a formula $\alpha\beta\gamma\delta$ set inline so
that the math list machinery runs inside a paragraph. A displayed formula
follows:
\[ \lim_{n\to\infty}\left(1+\frac{1}{n}\right)^{n}=e . \]
And a little more prose to fill the page out, with an em-dash --- and
quotation marks ``like these'' and a footnote.\footnote{A footnote, which is
an insertion, so the page builder must split it.}
\subsection{A subsection}
Running text again, with \textit{italics}, \textbf{bold}, and
\texttt{typewriter}, and a formula $\alpha\beta\gamma\delta$ set inline so
that the math list machinery runs inside a paragraph. A displayed formula
follows:
\[ \lim_{n\to\infty}\left(1+\frac{1}{n}\right)^{n}=e . \]
And a little more prose to fill the page out, with an em-dash --- and
quotation marks ``like these'' and a footnote.\footnote{A footnote, which is
an insertion, so the page builder must split it.}
\subsection{A subsection}
Running text again, with \textit{italics}, \textbf{bold}, and
\texttt{typewriter}, and a formula $\alpha\beta\gamma\delta$ set inline so
that the math list machinery runs inside a paragraph. A displayed formula
follows:
\[ \lim_{n\to\infty}\left(1+\frac{1}{n}\right)^{n}=e . \]
And a little more prose to fill the page out, with an em-dash --- and
quotation marks ``like these'' and a footnote.\footnote{A footnote, which is
an insertion, so the page builder must split it.}
\section{Prose that fills lines and breaks pages}
\label{sec:prose8-2}
Ordinary running text is what a typesetting engine spends most of its effort
upon, so the training document supplies a great deal of it. The words are
chosen for their length and their hyphenation, with incomprehensibility and
antidisestablishmentarianism thrown in so that the patterns are exercised.
Cross references such as \cref{sec:prose8-2} and citations of a page such
as \pageref{sec:prose8-2} make the auxiliary machinery run, and a
paragraph long enough to break over several lines gives the line breaker
something to weigh.
\begin{proposition}\label{prop:8-2}
For every $n\in\mathbb{N}$ the sum $\sum_{k=1}^{n}k$ equals
$\tfrac{1}{2}n(n+1)$.
\end{proposition}
\begin{proof}
Induction on $n$; the base case is $\pairs{1}{1}$ and the step adds $n+1$ to
both sides of the identity, which is what \cref{prop:8-2} asserts.
\end{proof}
\begin{equation}\label{eq:8-2}
  \int_{0}^{\infty}\frac{\vecn{x}^{s-1}}{e^{\vecn{x}}-1}\,d\vecn{x}
    = \Gamma(s)\zeta(s),\qquad \Re(s)>1 .
\end{equation}
\begin{align}
  \mathcal{L}[f](s) &= \int_{0}^{\infty} e^{-st} f(t)\,dt, \\
  \widehat{f}(\xi)  &= \int_{-\infty}^{\infty} f(x) e^{-2\pi i x\xi}\,dx .
\end{align}
\par\noindent\begin{center}
\begin{tabular}{lrrr}
\hline
Name & One & Two & Three \\
\hline
Alpha & 1 & 2 & 3 \\
Beta & 4 & 5 & 6 \\
Gamma & 7 & 8 & 9 \\
\hline
\end{tabular}
\end{center}
\begin{itemize}
  \item A first item with a formula $a^2+b^2=c^2$ in it.
  \item A second item that refers to \eqref{eq:8-2}.
  \item A third item, longer, so that the list breaks across a line and the
    paragraph builder has something to weigh when it chooses the break.
\end{itemize}
\subsection{A subsection}
Running text again, with \textit{italics}, \textbf{bold}, and
\texttt{typewriter}, and a formula $\alpha\beta\gamma\delta$ set inline so
that the math list machinery runs inside a paragraph. A displayed formula
follows:
\[ \lim_{n\to\infty}\left(1+\frac{1}{n}\right)^{n}=e . \]
And a little more prose to fill the page out, with an em-dash --- and
quotation marks ``like these'' and a footnote.\footnote{A footnote, which is
an insertion, so the page builder must split it.}
\subsection{A subsection}
Running text again, with \textit{italics}, \textbf{bold}, and
\texttt{typewriter}, and a formula $\alpha\beta\gamma\delta$ set inline so
that the math list machinery runs inside a paragraph. A displayed formula
follows:
\[ \lim_{n\to\infty}\left(1+\frac{1}{n}\right)^{n}=e . \]
And a little more prose to fill the page out, with an em-dash --- and
quotation marks ``like these'' and a footnote.\footnote{A footnote, which is
an insertion, so the page builder must split it.}
\subsection{A subsection}
Running text again, with \textit{italics}, \textbf{bold}, and
\texttt{typewriter}, and a formula $\alpha\beta\gamma\delta$ set inline so
that the math list machinery runs inside a paragraph. A displayed formula
follows:
\[ \lim_{n\to\infty}\left(1+\frac{1}{n}\right)^{n}=e . \]
And a little more prose to fill the page out, with an em-dash --- and
quotation marks ``like these'' and a footnote.\footnote{A footnote, which is
an insertion, so the page builder must split it.}
\section{Prose that fills lines and breaks pages}
\label{sec:prose8-3}
Ordinary running text is what a typesetting engine spends most of its effort
upon, so the training document supplies a great deal of it. The words are
chosen for their length and their hyphenation, with incomprehensibility and
antidisestablishmentarianism thrown in so that the patterns are exercised.
Cross references such as \cref{sec:prose8-3} and citations of a page such
as \pageref{sec:prose8-3} make the auxiliary machinery run, and a
paragraph long enough to break over several lines gives the line breaker
something to weigh.
\begin{proposition}\label{prop:8-3}
For every $n\in\mathbb{N}$ the sum $\sum_{k=1}^{n}k$ equals
$\tfrac{1}{2}n(n+1)$.
\end{proposition}
\begin{proof}
Induction on $n$; the base case is $\pairs{1}{1}$ and the step adds $n+1$ to
both sides of the identity, which is what \cref{prop:8-3} asserts.
\end{proof}
\begin{equation}\label{eq:8-3}
  \int_{0}^{\infty}\frac{\vecn{x}^{s-1}}{e^{\vecn{x}}-1}\,d\vecn{x}
    = \Gamma(s)\zeta(s),\qquad \Re(s)>1 .
\end{equation}
\begin{align}
  \mathcal{L}[f](s) &= \int_{0}^{\infty} e^{-st} f(t)\,dt, \\
  \widehat{f}(\xi)  &= \int_{-\infty}^{\infty} f(x) e^{-2\pi i x\xi}\,dx .
\end{align}
\par\noindent\begin{center}
\begin{tabular}{lrrr}
\hline
Name & One & Two & Three \\
\hline
Alpha & 1 & 2 & 3 \\
Beta & 4 & 5 & 6 \\
Gamma & 7 & 8 & 9 \\
\hline
\end{tabular}
\end{center}
\begin{itemize}
  \item A first item with a formula $a^2+b^2=c^2$ in it.
  \item A second item that refers to \eqref{eq:8-3}.
  \item A third item, longer, so that the list breaks across a line and the
    paragraph builder has something to weigh when it chooses the break.
\end{itemize}
\subsection{A subsection}
Running text again, with \textit{italics}, \textbf{bold}, and
\texttt{typewriter}, and a formula $\alpha\beta\gamma\delta$ set inline so
that the math list machinery runs inside a paragraph. A displayed formula
follows:
\[ \lim_{n\to\infty}\left(1+\frac{1}{n}\right)^{n}=e . \]
And a little more prose to fill the page out, with an em-dash --- and
quotation marks ``like these'' and a footnote.\footnote{A footnote, which is
an insertion, so the page builder must split it.}
\subsection{A subsection}
Running text again, with \textit{italics}, \textbf{bold}, and
\texttt{typewriter}, and a formula $\alpha\beta\gamma\delta$ set inline so
that the math list machinery runs inside a paragraph. A displayed formula
follows:
\[ \lim_{n\to\infty}\left(1+\frac{1}{n}\right)^{n}=e . \]
And a little more prose to fill the page out, with an em-dash --- and
quotation marks ``like these'' and a footnote.\footnote{A footnote, which is
an insertion, so the page builder must split it.}
\subsection{A subsection}
Running text again, with \textit{italics}, \textbf{bold}, and
\texttt{typewriter}, and a formula $\alpha\beta\gamma\delta$ set inline so
that the math list machinery runs inside a paragraph. A displayed formula
follows:
\[ \lim_{n\to\infty}\left(1+\frac{1}{n}\right)^{n}=e . \]
And a little more prose to fill the page out, with an em-dash --- and
quotation marks ``like these'' and a footnote.\footnote{A footnote, which is
an insertion, so the page builder must split it.}
\section{Prose that fills lines and breaks pages}
\label{sec:prose8-4}
Ordinary running text is what a typesetting engine spends most of its effort
upon, so the training document supplies a great deal of it. The words are
chosen for their length and their hyphenation, with incomprehensibility and
antidisestablishmentarianism thrown in so that the patterns are exercised.
Cross references such as \cref{sec:prose8-4} and citations of a page such
as \pageref{sec:prose8-4} make the auxiliary machinery run, and a
paragraph long enough to break over several lines gives the line breaker
something to weigh.
\begin{proposition}\label{prop:8-4}
For every $n\in\mathbb{N}$ the sum $\sum_{k=1}^{n}k$ equals
$\tfrac{1}{2}n(n+1)$.
\end{proposition}
\begin{proof}
Induction on $n$; the base case is $\pairs{1}{1}$ and the step adds $n+1$ to
both sides of the identity, which is what \cref{prop:8-4} asserts.
\end{proof}
\begin{equation}\label{eq:8-4}
  \int_{0}^{\infty}\frac{\vecn{x}^{s-1}}{e^{\vecn{x}}-1}\,d\vecn{x}
    = \Gamma(s)\zeta(s),\qquad \Re(s)>1 .
\end{equation}
\begin{align}
  \mathcal{L}[f](s) &= \int_{0}^{\infty} e^{-st} f(t)\,dt, \\
  \widehat{f}(\xi)  &= \int_{-\infty}^{\infty} f(x) e^{-2\pi i x\xi}\,dx .
\end{align}
\par\noindent\begin{center}
\begin{tabular}{lrrr}
\hline
Name & One & Two & Three \\
\hline
Alpha & 1 & 2 & 3 \\
Beta & 4 & 5 & 6 \\
Gamma & 7 & 8 & 9 \\
\hline
\end{tabular}
\end{center}
\begin{itemize}
  \item A first item with a formula $a^2+b^2=c^2$ in it.
  \item A second item that refers to \eqref{eq:8-4}.
  \item A third item, longer, so that the list breaks across a line and the
    paragraph builder has something to weigh when it chooses the break.
\end{itemize}
\subsection{A subsection}
Running text again, with \textit{italics}, \textbf{bold}, and
\texttt{typewriter}, and a formula $\alpha\beta\gamma\delta$ set inline so
that the math list machinery runs inside a paragraph. A displayed formula
follows:
\[ \lim_{n\to\infty}\left(1+\frac{1}{n}\right)^{n}=e . \]
And a little more prose to fill the page out, with an em-dash --- and
quotation marks ``like these'' and a footnote.\footnote{A footnote, which is
an insertion, so the page builder must split it.}
\subsection{A subsection}
Running text again, with \textit{italics}, \textbf{bold}, and
\texttt{typewriter}, and a formula $\alpha\beta\gamma\delta$ set inline so
that the math list machinery runs inside a paragraph. A displayed formula
follows:
\[ \lim_{n\to\infty}\left(1+\frac{1}{n}\right)^{n}=e . \]
And a little more prose to fill the page out, with an em-dash --- and
quotation marks ``like these'' and a footnote.\footnote{A footnote, which is
an insertion, so the page builder must split it.}
\subsection{A subsection}
Running text again, with \textit{italics}, \textbf{bold}, and
\texttt{typewriter}, and a formula $\alpha\beta\gamma\delta$ set inline so
that the math list machinery runs inside a paragraph. A displayed formula
follows:
\[ \lim_{n\to\infty}\left(1+\frac{1}{n}\right)^{n}=e . \]
And a little more prose to fill the page out, with an em-dash --- and
quotation marks ``like these'' and a footnote.\footnote{A footnote, which is
an insertion, so the page builder must split it.}
\chapter{A chapter of ordinary matter, number 9}
\section{Prose that fills lines and breaks pages}
\label{sec:prose9-1}
Ordinary running text is what a typesetting engine spends most of its effort
upon, so the training document supplies a great deal of it. The words are
chosen for their length and their hyphenation, with incomprehensibility and
antidisestablishmentarianism thrown in so that the patterns are exercised.
Cross references such as \cref{sec:prose9-1} and citations of a page such
as \pageref{sec:prose9-1} make the auxiliary machinery run, and a
paragraph long enough to break over several lines gives the line breaker
something to weigh.
\begin{proposition}\label{prop:9-1}
For every $n\in\mathbb{N}$ the sum $\sum_{k=1}^{n}k$ equals
$\tfrac{1}{2}n(n+1)$.
\end{proposition}
\begin{proof}
Induction on $n$; the base case is $\pairs{1}{1}$ and the step adds $n+1$ to
both sides of the identity, which is what \cref{prop:9-1} asserts.
\end{proof}
\begin{equation}\label{eq:9-1}
  \int_{0}^{\infty}\frac{\vecn{x}^{s-1}}{e^{\vecn{x}}-1}\,d\vecn{x}
    = \Gamma(s)\zeta(s),\qquad \Re(s)>1 .
\end{equation}
\begin{align}
  \mathcal{L}[f](s) &= \int_{0}^{\infty} e^{-st} f(t)\,dt, \\
  \widehat{f}(\xi)  &= \int_{-\infty}^{\infty} f(x) e^{-2\pi i x\xi}\,dx .
\end{align}
\par\noindent\begin{center}
\begin{tabular}{lrrr}
\hline
Name & One & Two & Three \\
\hline
Alpha & 1 & 2 & 3 \\
Beta & 4 & 5 & 6 \\
Gamma & 7 & 8 & 9 \\
\hline
\end{tabular}
\end{center}
\begin{itemize}
  \item A first item with a formula $a^2+b^2=c^2$ in it.
  \item A second item that refers to \eqref{eq:9-1}.
  \item A third item, longer, so that the list breaks across a line and the
    paragraph builder has something to weigh when it chooses the break.
\end{itemize}
\subsection{A subsection}
Running text again, with \textit{italics}, \textbf{bold}, and
\texttt{typewriter}, and a formula $\alpha\beta\gamma\delta$ set inline so
that the math list machinery runs inside a paragraph. A displayed formula
follows:
\[ \lim_{n\to\infty}\left(1+\frac{1}{n}\right)^{n}=e . \]
And a little more prose to fill the page out, with an em-dash --- and
quotation marks ``like these'' and a footnote.\footnote{A footnote, which is
an insertion, so the page builder must split it.}
\subsection{A subsection}
Running text again, with \textit{italics}, \textbf{bold}, and
\texttt{typewriter}, and a formula $\alpha\beta\gamma\delta$ set inline so
that the math list machinery runs inside a paragraph. A displayed formula
follows:
\[ \lim_{n\to\infty}\left(1+\frac{1}{n}\right)^{n}=e . \]
And a little more prose to fill the page out, with an em-dash --- and
quotation marks ``like these'' and a footnote.\footnote{A footnote, which is
an insertion, so the page builder must split it.}
\subsection{A subsection}
Running text again, with \textit{italics}, \textbf{bold}, and
\texttt{typewriter}, and a formula $\alpha\beta\gamma\delta$ set inline so
that the math list machinery runs inside a paragraph. A displayed formula
follows:
\[ \lim_{n\to\infty}\left(1+\frac{1}{n}\right)^{n}=e . \]
And a little more prose to fill the page out, with an em-dash --- and
quotation marks ``like these'' and a footnote.\footnote{A footnote, which is
an insertion, so the page builder must split it.}
\section{Prose that fills lines and breaks pages}
\label{sec:prose9-2}
Ordinary running text is what a typesetting engine spends most of its effort
upon, so the training document supplies a great deal of it. The words are
chosen for their length and their hyphenation, with incomprehensibility and
antidisestablishmentarianism thrown in so that the patterns are exercised.
Cross references such as \cref{sec:prose9-2} and citations of a page such
as \pageref{sec:prose9-2} make the auxiliary machinery run, and a
paragraph long enough to break over several lines gives the line breaker
something to weigh.
\begin{proposition}\label{prop:9-2}
For every $n\in\mathbb{N}$ the sum $\sum_{k=1}^{n}k$ equals
$\tfrac{1}{2}n(n+1)$.
\end{proposition}
\begin{proof}
Induction on $n$; the base case is $\pairs{1}{1}$ and the step adds $n+1$ to
both sides of the identity, which is what \cref{prop:9-2} asserts.
\end{proof}
\begin{equation}\label{eq:9-2}
  \int_{0}^{\infty}\frac{\vecn{x}^{s-1}}{e^{\vecn{x}}-1}\,d\vecn{x}
    = \Gamma(s)\zeta(s),\qquad \Re(s)>1 .
\end{equation}
\begin{align}
  \mathcal{L}[f](s) &= \int_{0}^{\infty} e^{-st} f(t)\,dt, \\
  \widehat{f}(\xi)  &= \int_{-\infty}^{\infty} f(x) e^{-2\pi i x\xi}\,dx .
\end{align}
\par\noindent\begin{center}
\begin{tabular}{lrrr}
\hline
Name & One & Two & Three \\
\hline
Alpha & 1 & 2 & 3 \\
Beta & 4 & 5 & 6 \\
Gamma & 7 & 8 & 9 \\
\hline
\end{tabular}
\end{center}
\begin{itemize}
  \item A first item with a formula $a^2+b^2=c^2$ in it.
  \item A second item that refers to \eqref{eq:9-2}.
  \item A third item, longer, so that the list breaks across a line and the
    paragraph builder has something to weigh when it chooses the break.
\end{itemize}
\subsection{A subsection}
Running text again, with \textit{italics}, \textbf{bold}, and
\texttt{typewriter}, and a formula $\alpha\beta\gamma\delta$ set inline so
that the math list machinery runs inside a paragraph. A displayed formula
follows:
\[ \lim_{n\to\infty}\left(1+\frac{1}{n}\right)^{n}=e . \]
And a little more prose to fill the page out, with an em-dash --- and
quotation marks ``like these'' and a footnote.\footnote{A footnote, which is
an insertion, so the page builder must split it.}
\subsection{A subsection}
Running text again, with \textit{italics}, \textbf{bold}, and
\texttt{typewriter}, and a formula $\alpha\beta\gamma\delta$ set inline so
that the math list machinery runs inside a paragraph. A displayed formula
follows:
\[ \lim_{n\to\infty}\left(1+\frac{1}{n}\right)^{n}=e . \]
And a little more prose to fill the page out, with an em-dash --- and
quotation marks ``like these'' and a footnote.\footnote{A footnote, which is
an insertion, so the page builder must split it.}
\subsection{A subsection}
Running text again, with \textit{italics}, \textbf{bold}, and
\texttt{typewriter}, and a formula $\alpha\beta\gamma\delta$ set inline so
that the math list machinery runs inside a paragraph. A displayed formula
follows:
\[ \lim_{n\to\infty}\left(1+\frac{1}{n}\right)^{n}=e . \]
And a little more prose to fill the page out, with an em-dash --- and
quotation marks ``like these'' and a footnote.\footnote{A footnote, which is
an insertion, so the page builder must split it.}
\section{Prose that fills lines and breaks pages}
\label{sec:prose9-3}
Ordinary running text is what a typesetting engine spends most of its effort
upon, so the training document supplies a great deal of it. The words are
chosen for their length and their hyphenation, with incomprehensibility and
antidisestablishmentarianism thrown in so that the patterns are exercised.
Cross references such as \cref{sec:prose9-3} and citations of a page such
as \pageref{sec:prose9-3} make the auxiliary machinery run, and a
paragraph long enough to break over several lines gives the line breaker
something to weigh.
\begin{proposition}\label{prop:9-3}
For every $n\in\mathbb{N}$ the sum $\sum_{k=1}^{n}k$ equals
$\tfrac{1}{2}n(n+1)$.
\end{proposition}
\begin{proof}
Induction on $n$; the base case is $\pairs{1}{1}$ and the step adds $n+1$ to
both sides of the identity, which is what \cref{prop:9-3} asserts.
\end{proof}
\begin{equation}\label{eq:9-3}
  \int_{0}^{\infty}\frac{\vecn{x}^{s-1}}{e^{\vecn{x}}-1}\,d\vecn{x}
    = \Gamma(s)\zeta(s),\qquad \Re(s)>1 .
\end{equation}
\begin{align}
  \mathcal{L}[f](s) &= \int_{0}^{\infty} e^{-st} f(t)\,dt, \\
  \widehat{f}(\xi)  &= \int_{-\infty}^{\infty} f(x) e^{-2\pi i x\xi}\,dx .
\end{align}
\par\noindent\begin{center}
\begin{tabular}{lrrr}
\hline
Name & One & Two & Three \\
\hline
Alpha & 1 & 2 & 3 \\
Beta & 4 & 5 & 6 \\
Gamma & 7 & 8 & 9 \\
\hline
\end{tabular}
\end{center}
\begin{itemize}
  \item A first item with a formula $a^2+b^2=c^2$ in it.
  \item A second item that refers to \eqref{eq:9-3}.
  \item A third item, longer, so that the list breaks across a line and the
    paragraph builder has something to weigh when it chooses the break.
\end{itemize}
\subsection{A subsection}
Running text again, with \textit{italics}, \textbf{bold}, and
\texttt{typewriter}, and a formula $\alpha\beta\gamma\delta$ set inline so
that the math list machinery runs inside a paragraph. A displayed formula
follows:
\[ \lim_{n\to\infty}\left(1+\frac{1}{n}\right)^{n}=e . \]
And a little more prose to fill the page out, with an em-dash --- and
quotation marks ``like these'' and a footnote.\footnote{A footnote, which is
an insertion, so the page builder must split it.}
\subsection{A subsection}
Running text again, with \textit{italics}, \textbf{bold}, and
\texttt{typewriter}, and a formula $\alpha\beta\gamma\delta$ set inline so
that the math list machinery runs inside a paragraph. A displayed formula
follows:
\[ \lim_{n\to\infty}\left(1+\frac{1}{n}\right)^{n}=e . \]
And a little more prose to fill the page out, with an em-dash --- and
quotation marks ``like these'' and a footnote.\footnote{A footnote, which is
an insertion, so the page builder must split it.}
\subsection{A subsection}
Running text again, with \textit{italics}, \textbf{bold}, and
\texttt{typewriter}, and a formula $\alpha\beta\gamma\delta$ set inline so
that the math list machinery runs inside a paragraph. A displayed formula
follows:
\[ \lim_{n\to\infty}\left(1+\frac{1}{n}\right)^{n}=e . \]
And a little more prose to fill the page out, with an em-dash --- and
quotation marks ``like these'' and a footnote.\footnote{A footnote, which is
an insertion, so the page builder must split it.}
\section{Prose that fills lines and breaks pages}
\label{sec:prose9-4}
Ordinary running text is what a typesetting engine spends most of its effort
upon, so the training document supplies a great deal of it. The words are
chosen for their length and their hyphenation, with incomprehensibility and
antidisestablishmentarianism thrown in so that the patterns are exercised.
Cross references such as \cref{sec:prose9-4} and citations of a page such
as \pageref{sec:prose9-4} make the auxiliary machinery run, and a
paragraph long enough to break over several lines gives the line breaker
something to weigh.
\begin{proposition}\label{prop:9-4}
For every $n\in\mathbb{N}$ the sum $\sum_{k=1}^{n}k$ equals
$\tfrac{1}{2}n(n+1)$.
\end{proposition}
\begin{proof}
Induction on $n$; the base case is $\pairs{1}{1}$ and the step adds $n+1$ to
both sides of the identity, which is what \cref{prop:9-4} asserts.
\end{proof}
\begin{equation}\label{eq:9-4}
  \int_{0}^{\infty}\frac{\vecn{x}^{s-1}}{e^{\vecn{x}}-1}\,d\vecn{x}
    = \Gamma(s)\zeta(s),\qquad \Re(s)>1 .
\end{equation}
\begin{align}
  \mathcal{L}[f](s) &= \int_{0}^{\infty} e^{-st} f(t)\,dt, \\
  \widehat{f}(\xi)  &= \int_{-\infty}^{\infty} f(x) e^{-2\pi i x\xi}\,dx .
\end{align}
\par\noindent\begin{center}
\begin{tabular}{lrrr}
\hline
Name & One & Two & Three \\
\hline
Alpha & 1 & 2 & 3 \\
Beta & 4 & 5 & 6 \\
Gamma & 7 & 8 & 9 \\
\hline
\end{tabular}
\end{center}
\begin{itemize}
  \item A first item with a formula $a^2+b^2=c^2$ in it.
  \item A second item that refers to \eqref{eq:9-4}.
  \item A third item, longer, so that the list breaks across a line and the
    paragraph builder has something to weigh when it chooses the break.
\end{itemize}
\subsection{A subsection}
Running text again, with \textit{italics}, \textbf{bold}, and
\texttt{typewriter}, and a formula $\alpha\beta\gamma\delta$ set inline so
that the math list machinery runs inside a paragraph. A displayed formula
follows:
\[ \lim_{n\to\infty}\left(1+\frac{1}{n}\right)^{n}=e . \]
And a little more prose to fill the page out, with an em-dash --- and
quotation marks ``like these'' and a footnote.\footnote{A footnote, which is
an insertion, so the page builder must split it.}
\subsection{A subsection}
Running text again, with \textit{italics}, \textbf{bold}, and
\texttt{typewriter}, and a formula $\alpha\beta\gamma\delta$ set inline so
that the math list machinery runs inside a paragraph. A displayed formula
follows:
\[ \lim_{n\to\infty}\left(1+\frac{1}{n}\right)^{n}=e . \]
And a little more prose to fill the page out, with an em-dash --- and
quotation marks ``like these'' and a footnote.\footnote{A footnote, which is
an insertion, so the page builder must split it.}
\subsection{A subsection}
Running text again, with \textit{italics}, \textbf{bold}, and
\texttt{typewriter}, and a formula $\alpha\beta\gamma\delta$ set inline so
that the math list machinery runs inside a paragraph. A displayed formula
follows:
\[ \lim_{n\to\infty}\left(1+\frac{1}{n}\right)^{n}=e . \]
And a little more prose to fill the page out, with an em-dash --- and
quotation marks ``like these'' and a footnote.\footnote{A footnote, which is
an insertion, so the page builder must split it.}
\chapter{A chapter of ordinary matter, number 10}
\section{Prose that fills lines and breaks pages}
\label{sec:prose10-1}
Ordinary running text is what a typesetting engine spends most of its effort
upon, so the training document supplies a great deal of it. The words are
chosen for their length and their hyphenation, with incomprehensibility and
antidisestablishmentarianism thrown in so that the patterns are exercised.
Cross references such as \cref{sec:prose10-1} and citations of a page such
as \pageref{sec:prose10-1} make the auxiliary machinery run, and a
paragraph long enough to break over several lines gives the line breaker
something to weigh.
\begin{proposition}\label{prop:10-1}
For every $n\in\mathbb{N}$ the sum $\sum_{k=1}^{n}k$ equals
$\tfrac{1}{2}n(n+1)$.
\end{proposition}
\begin{proof}
Induction on $n$; the base case is $\pairs{1}{1}$ and the step adds $n+1$ to
both sides of the identity, which is what \cref{prop:10-1} asserts.
\end{proof}
\begin{equation}\label{eq:10-1}
  \int_{0}^{\infty}\frac{\vecn{x}^{s-1}}{e^{\vecn{x}}-1}\,d\vecn{x}
    = \Gamma(s)\zeta(s),\qquad \Re(s)>1 .
\end{equation}
\begin{align}
  \mathcal{L}[f](s) &= \int_{0}^{\infty} e^{-st} f(t)\,dt, \\
  \widehat{f}(\xi)  &= \int_{-\infty}^{\infty} f(x) e^{-2\pi i x\xi}\,dx .
\end{align}
\par\noindent\begin{center}
\begin{tabular}{lrrr}
\hline
Name & One & Two & Three \\
\hline
Alpha & 1 & 2 & 3 \\
Beta & 4 & 5 & 6 \\
Gamma & 7 & 8 & 9 \\
\hline
\end{tabular}
\end{center}
\begin{itemize}
  \item A first item with a formula $a^2+b^2=c^2$ in it.
  \item A second item that refers to \eqref{eq:10-1}.
  \item A third item, longer, so that the list breaks across a line and the
    paragraph builder has something to weigh when it chooses the break.
\end{itemize}
\subsection{A subsection}
Running text again, with \textit{italics}, \textbf{bold}, and
\texttt{typewriter}, and a formula $\alpha\beta\gamma\delta$ set inline so
that the math list machinery runs inside a paragraph. A displayed formula
follows:
\[ \lim_{n\to\infty}\left(1+\frac{1}{n}\right)^{n}=e . \]
And a little more prose to fill the page out, with an em-dash --- and
quotation marks ``like these'' and a footnote.\footnote{A footnote, which is
an insertion, so the page builder must split it.}
\subsection{A subsection}
Running text again, with \textit{italics}, \textbf{bold}, and
\texttt{typewriter}, and a formula $\alpha\beta\gamma\delta$ set inline so
that the math list machinery runs inside a paragraph. A displayed formula
follows:
\[ \lim_{n\to\infty}\left(1+\frac{1}{n}\right)^{n}=e . \]
And a little more prose to fill the page out, with an em-dash --- and
quotation marks ``like these'' and a footnote.\footnote{A footnote, which is
an insertion, so the page builder must split it.}
\subsection{A subsection}
Running text again, with \textit{italics}, \textbf{bold}, and
\texttt{typewriter}, and a formula $\alpha\beta\gamma\delta$ set inline so
that the math list machinery runs inside a paragraph. A displayed formula
follows:
\[ \lim_{n\to\infty}\left(1+\frac{1}{n}\right)^{n}=e . \]
And a little more prose to fill the page out, with an em-dash --- and
quotation marks ``like these'' and a footnote.\footnote{A footnote, which is
an insertion, so the page builder must split it.}
\section{Prose that fills lines and breaks pages}
\label{sec:prose10-2}
Ordinary running text is what a typesetting engine spends most of its effort
upon, so the training document supplies a great deal of it. The words are
chosen for their length and their hyphenation, with incomprehensibility and
antidisestablishmentarianism thrown in so that the patterns are exercised.
Cross references such as \cref{sec:prose10-2} and citations of a page such
as \pageref{sec:prose10-2} make the auxiliary machinery run, and a
paragraph long enough to break over several lines gives the line breaker
something to weigh.
\begin{proposition}\label{prop:10-2}
For every $n\in\mathbb{N}$ the sum $\sum_{k=1}^{n}k$ equals
$\tfrac{1}{2}n(n+1)$.
\end{proposition}
\begin{proof}
Induction on $n$; the base case is $\pairs{1}{1}$ and the step adds $n+1$ to
both sides of the identity, which is what \cref{prop:10-2} asserts.
\end{proof}
\begin{equation}\label{eq:10-2}
  \int_{0}^{\infty}\frac{\vecn{x}^{s-1}}{e^{\vecn{x}}-1}\,d\vecn{x}
    = \Gamma(s)\zeta(s),\qquad \Re(s)>1 .
\end{equation}
\begin{align}
  \mathcal{L}[f](s) &= \int_{0}^{\infty} e^{-st} f(t)\,dt, \\
  \widehat{f}(\xi)  &= \int_{-\infty}^{\infty} f(x) e^{-2\pi i x\xi}\,dx .
\end{align}
\par\noindent\begin{center}
\begin{tabular}{lrrr}
\hline
Name & One & Two & Three \\
\hline
Alpha & 1 & 2 & 3 \\
Beta & 4 & 5 & 6 \\
Gamma & 7 & 8 & 9 \\
\hline
\end{tabular}
\end{center}
\begin{itemize}
  \item A first item with a formula $a^2+b^2=c^2$ in it.
  \item A second item that refers to \eqref{eq:10-2}.
  \item A third item, longer, so that the list breaks across a line and the
    paragraph builder has something to weigh when it chooses the break.
\end{itemize}
\subsection{A subsection}
Running text again, with \textit{italics}, \textbf{bold}, and
\texttt{typewriter}, and a formula $\alpha\beta\gamma\delta$ set inline so
that the math list machinery runs inside a paragraph. A displayed formula
follows:
\[ \lim_{n\to\infty}\left(1+\frac{1}{n}\right)^{n}=e . \]
And a little more prose to fill the page out, with an em-dash --- and
quotation marks ``like these'' and a footnote.\footnote{A footnote, which is
an insertion, so the page builder must split it.}
\subsection{A subsection}
Running text again, with \textit{italics}, \textbf{bold}, and
\texttt{typewriter}, and a formula $\alpha\beta\gamma\delta$ set inline so
that the math list machinery runs inside a paragraph. A displayed formula
follows:
\[ \lim_{n\to\infty}\left(1+\frac{1}{n}\right)^{n}=e . \]
And a little more prose to fill the page out, with an em-dash --- and
quotation marks ``like these'' and a footnote.\footnote{A footnote, which is
an insertion, so the page builder must split it.}
\subsection{A subsection}
Running text again, with \textit{italics}, \textbf{bold}, and
\texttt{typewriter}, and a formula $\alpha\beta\gamma\delta$ set inline so
that the math list machinery runs inside a paragraph. A displayed formula
follows:
\[ \lim_{n\to\infty}\left(1+\frac{1}{n}\right)^{n}=e . \]
And a little more prose to fill the page out, with an em-dash --- and
quotation marks ``like these'' and a footnote.\footnote{A footnote, which is
an insertion, so the page builder must split it.}
\section{Prose that fills lines and breaks pages}
\label{sec:prose10-3}
Ordinary running text is what a typesetting engine spends most of its effort
upon, so the training document supplies a great deal of it. The words are
chosen for their length and their hyphenation, with incomprehensibility and
antidisestablishmentarianism thrown in so that the patterns are exercised.
Cross references such as \cref{sec:prose10-3} and citations of a page such
as \pageref{sec:prose10-3} make the auxiliary machinery run, and a
paragraph long enough to break over several lines gives the line breaker
something to weigh.
\begin{proposition}\label{prop:10-3}
For every $n\in\mathbb{N}$ the sum $\sum_{k=1}^{n}k$ equals
$\tfrac{1}{2}n(n+1)$.
\end{proposition}
\begin{proof}
Induction on $n$; the base case is $\pairs{1}{1}$ and the step adds $n+1$ to
both sides of the identity, which is what \cref{prop:10-3} asserts.
\end{proof}
\begin{equation}\label{eq:10-3}
  \int_{0}^{\infty}\frac{\vecn{x}^{s-1}}{e^{\vecn{x}}-1}\,d\vecn{x}
    = \Gamma(s)\zeta(s),\qquad \Re(s)>1 .
\end{equation}
\begin{align}
  \mathcal{L}[f](s) &= \int_{0}^{\infty} e^{-st} f(t)\,dt, \\
  \widehat{f}(\xi)  &= \int_{-\infty}^{\infty} f(x) e^{-2\pi i x\xi}\,dx .
\end{align}
\par\noindent\begin{center}
\begin{tabular}{lrrr}
\hline
Name & One & Two & Three \\
\hline
Alpha & 1 & 2 & 3 \\
Beta & 4 & 5 & 6 \\
Gamma & 7 & 8 & 9 \\
\hline
\end{tabular}
\end{center}
\begin{itemize}
  \item A first item with a formula $a^2+b^2=c^2$ in it.
  \item A second item that refers to \eqref{eq:10-3}.
  \item A third item, longer, so that the list breaks across a line and the
    paragraph builder has something to weigh when it chooses the break.
\end{itemize}
\subsection{A subsection}
Running text again, with \textit{italics}, \textbf{bold}, and
\texttt{typewriter}, and a formula $\alpha\beta\gamma\delta$ set inline so
that the math list machinery runs inside a paragraph. A displayed formula
follows:
\[ \lim_{n\to\infty}\left(1+\frac{1}{n}\right)^{n}=e . \]
And a little more prose to fill the page out, with an em-dash --- and
quotation marks ``like these'' and a footnote.\footnote{A footnote, which is
an insertion, so the page builder must split it.}
\subsection{A subsection}
Running text again, with \textit{italics}, \textbf{bold}, and
\texttt{typewriter}, and a formula $\alpha\beta\gamma\delta$ set inline so
that the math list machinery runs inside a paragraph. A displayed formula
follows:
\[ \lim_{n\to\infty}\left(1+\frac{1}{n}\right)^{n}=e . \]
And a little more prose to fill the page out, with an em-dash --- and
quotation marks ``like these'' and a footnote.\footnote{A footnote, which is
an insertion, so the page builder must split it.}
\subsection{A subsection}
Running text again, with \textit{italics}, \textbf{bold}, and
\texttt{typewriter}, and a formula $\alpha\beta\gamma\delta$ set inline so
that the math list machinery runs inside a paragraph. A displayed formula
follows:
\[ \lim_{n\to\infty}\left(1+\frac{1}{n}\right)^{n}=e . \]
And a little more prose to fill the page out, with an em-dash --- and
quotation marks ``like these'' and a footnote.\footnote{A footnote, which is
an insertion, so the page builder must split it.}
\section{Prose that fills lines and breaks pages}
\label{sec:prose10-4}
Ordinary running text is what a typesetting engine spends most of its effort
upon, so the training document supplies a great deal of it. The words are
chosen for their length and their hyphenation, with incomprehensibility and
antidisestablishmentarianism thrown in so that the patterns are exercised.
Cross references such as \cref{sec:prose10-4} and citations of a page such
as \pageref{sec:prose10-4} make the auxiliary machinery run, and a
paragraph long enough to break over several lines gives the line breaker
something to weigh.
\begin{proposition}\label{prop:10-4}
For every $n\in\mathbb{N}$ the sum $\sum_{k=1}^{n}k$ equals
$\tfrac{1}{2}n(n+1)$.
\end{proposition}
\begin{proof}
Induction on $n$; the base case is $\pairs{1}{1}$ and the step adds $n+1$ to
both sides of the identity, which is what \cref{prop:10-4} asserts.
\end{proof}
\begin{equation}\label{eq:10-4}
  \int_{0}^{\infty}\frac{\vecn{x}^{s-1}}{e^{\vecn{x}}-1}\,d\vecn{x}
    = \Gamma(s)\zeta(s),\qquad \Re(s)>1 .
\end{equation}
\begin{align}
  \mathcal{L}[f](s) &= \int_{0}^{\infty} e^{-st} f(t)\,dt, \\
  \widehat{f}(\xi)  &= \int_{-\infty}^{\infty} f(x) e^{-2\pi i x\xi}\,dx .
\end{align}
\par\noindent\begin{center}
\begin{tabular}{lrrr}
\hline
Name & One & Two & Three \\
\hline
Alpha & 1 & 2 & 3 \\
Beta & 4 & 5 & 6 \\
Gamma & 7 & 8 & 9 \\
\hline
\end{tabular}
\end{center}
\begin{itemize}
  \item A first item with a formula $a^2+b^2=c^2$ in it.
  \item A second item that refers to \eqref{eq:10-4}.
  \item A third item, longer, so that the list breaks across a line and the
    paragraph builder has something to weigh when it chooses the break.
\end{itemize}
\subsection{A subsection}
Running text again, with \textit{italics}, \textbf{bold}, and
\texttt{typewriter}, and a formula $\alpha\beta\gamma\delta$ set inline so
that the math list machinery runs inside a paragraph. A displayed formula
follows:
\[ \lim_{n\to\infty}\left(1+\frac{1}{n}\right)^{n}=e . \]
And a little more prose to fill the page out, with an em-dash --- and
quotation marks ``like these'' and a footnote.\footnote{A footnote, which is
an insertion, so the page builder must split it.}
\subsection{A subsection}
Running text again, with \textit{italics}, \textbf{bold}, and
\texttt{typewriter}, and a formula $\alpha\beta\gamma\delta$ set inline so
that the math list machinery runs inside a paragraph. A displayed formula
follows:
\[ \lim_{n\to\infty}\left(1+\frac{1}{n}\right)^{n}=e . \]
And a little more prose to fill the page out, with an em-dash --- and
quotation marks ``like these'' and a footnote.\footnote{A footnote, which is
an insertion, so the page builder must split it.}
\subsection{A subsection}
Running text again, with \textit{italics}, \textbf{bold}, and
\texttt{typewriter}, and a formula $\alpha\beta\gamma\delta$ set inline so
that the math list machinery runs inside a paragraph. A displayed formula
follows:
\[ \lim_{n\to\infty}\left(1+\frac{1}{n}\right)^{n}=e . \]
And a little more prose to fill the page out, with an em-dash --- and
quotation marks ``like these'' and a footnote.\footnote{A footnote, which is
an insertion, so the page builder must split it.}
\chapter{A chapter of ordinary matter, number 11}
\section{Prose that fills lines and breaks pages}
\label{sec:prose11-1}
Ordinary running text is what a typesetting engine spends most of its effort
upon, so the training document supplies a great deal of it. The words are
chosen for their length and their hyphenation, with incomprehensibility and
antidisestablishmentarianism thrown in so that the patterns are exercised.
Cross references such as \cref{sec:prose11-1} and citations of a page such
as \pageref{sec:prose11-1} make the auxiliary machinery run, and a
paragraph long enough to break over several lines gives the line breaker
something to weigh.
\begin{proposition}\label{prop:11-1}
For every $n\in\mathbb{N}$ the sum $\sum_{k=1}^{n}k$ equals
$\tfrac{1}{2}n(n+1)$.
\end{proposition}
\begin{proof}
Induction on $n$; the base case is $\pairs{1}{1}$ and the step adds $n+1$ to
both sides of the identity, which is what \cref{prop:11-1} asserts.
\end{proof}
\begin{equation}\label{eq:11-1}
  \int_{0}^{\infty}\frac{\vecn{x}^{s-1}}{e^{\vecn{x}}-1}\,d\vecn{x}
    = \Gamma(s)\zeta(s),\qquad \Re(s)>1 .
\end{equation}
\begin{align}
  \mathcal{L}[f](s) &= \int_{0}^{\infty} e^{-st} f(t)\,dt, \\
  \widehat{f}(\xi)  &= \int_{-\infty}^{\infty} f(x) e^{-2\pi i x\xi}\,dx .
\end{align}
\par\noindent\begin{center}
\begin{tabular}{lrrr}
\hline
Name & One & Two & Three \\
\hline
Alpha & 1 & 2 & 3 \\
Beta & 4 & 5 & 6 \\
Gamma & 7 & 8 & 9 \\
\hline
\end{tabular}
\end{center}
\begin{itemize}
  \item A first item with a formula $a^2+b^2=c^2$ in it.
  \item A second item that refers to \eqref{eq:11-1}.
  \item A third item, longer, so that the list breaks across a line and the
    paragraph builder has something to weigh when it chooses the break.
\end{itemize}
\subsection{A subsection}
Running text again, with \textit{italics}, \textbf{bold}, and
\texttt{typewriter}, and a formula $\alpha\beta\gamma\delta$ set inline so
that the math list machinery runs inside a paragraph. A displayed formula
follows:
\[ \lim_{n\to\infty}\left(1+\frac{1}{n}\right)^{n}=e . \]
And a little more prose to fill the page out, with an em-dash --- and
quotation marks ``like these'' and a footnote.\footnote{A footnote, which is
an insertion, so the page builder must split it.}
\subsection{A subsection}
Running text again, with \textit{italics}, \textbf{bold}, and
\texttt{typewriter}, and a formula $\alpha\beta\gamma\delta$ set inline so
that the math list machinery runs inside a paragraph. A displayed formula
follows:
\[ \lim_{n\to\infty}\left(1+\frac{1}{n}\right)^{n}=e . \]
And a little more prose to fill the page out, with an em-dash --- and
quotation marks ``like these'' and a footnote.\footnote{A footnote, which is
an insertion, so the page builder must split it.}
\subsection{A subsection}
Running text again, with \textit{italics}, \textbf{bold}, and
\texttt{typewriter}, and a formula $\alpha\beta\gamma\delta$ set inline so
that the math list machinery runs inside a paragraph. A displayed formula
follows:
\[ \lim_{n\to\infty}\left(1+\frac{1}{n}\right)^{n}=e . \]
And a little more prose to fill the page out, with an em-dash --- and
quotation marks ``like these'' and a footnote.\footnote{A footnote, which is
an insertion, so the page builder must split it.}
\section{Prose that fills lines and breaks pages}
\label{sec:prose11-2}
Ordinary running text is what a typesetting engine spends most of its effort
upon, so the training document supplies a great deal of it. The words are
chosen for their length and their hyphenation, with incomprehensibility and
antidisestablishmentarianism thrown in so that the patterns are exercised.
Cross references such as \cref{sec:prose11-2} and citations of a page such
as \pageref{sec:prose11-2} make the auxiliary machinery run, and a
paragraph long enough to break over several lines gives the line breaker
something to weigh.
\begin{proposition}\label{prop:11-2}
For every $n\in\mathbb{N}$ the sum $\sum_{k=1}^{n}k$ equals
$\tfrac{1}{2}n(n+1)$.
\end{proposition}
\begin{proof}
Induction on $n$; the base case is $\pairs{1}{1}$ and the step adds $n+1$ to
both sides of the identity, which is what \cref{prop:11-2} asserts.
\end{proof}
\begin{equation}\label{eq:11-2}
  \int_{0}^{\infty}\frac{\vecn{x}^{s-1}}{e^{\vecn{x}}-1}\,d\vecn{x}
    = \Gamma(s)\zeta(s),\qquad \Re(s)>1 .
\end{equation}
\begin{align}
  \mathcal{L}[f](s) &= \int_{0}^{\infty} e^{-st} f(t)\,dt, \\
  \widehat{f}(\xi)  &= \int_{-\infty}^{\infty} f(x) e^{-2\pi i x\xi}\,dx .
\end{align}
\par\noindent\begin{center}
\begin{tabular}{lrrr}
\hline
Name & One & Two & Three \\
\hline
Alpha & 1 & 2 & 3 \\
Beta & 4 & 5 & 6 \\
Gamma & 7 & 8 & 9 \\
\hline
\end{tabular}
\end{center}
\begin{itemize}
  \item A first item with a formula $a^2+b^2=c^2$ in it.
  \item A second item that refers to \eqref{eq:11-2}.
  \item A third item, longer, so that the list breaks across a line and the
    paragraph builder has something to weigh when it chooses the break.
\end{itemize}
\subsection{A subsection}
Running text again, with \textit{italics}, \textbf{bold}, and
\texttt{typewriter}, and a formula $\alpha\beta\gamma\delta$ set inline so
that the math list machinery runs inside a paragraph. A displayed formula
follows:
\[ \lim_{n\to\infty}\left(1+\frac{1}{n}\right)^{n}=e . \]
And a little more prose to fill the page out, with an em-dash --- and
quotation marks ``like these'' and a footnote.\footnote{A footnote, which is
an insertion, so the page builder must split it.}
\subsection{A subsection}
Running text again, with \textit{italics}, \textbf{bold}, and
\texttt{typewriter}, and a formula $\alpha\beta\gamma\delta$ set inline so
that the math list machinery runs inside a paragraph. A displayed formula
follows:
\[ \lim_{n\to\infty}\left(1+\frac{1}{n}\right)^{n}=e . \]
And a little more prose to fill the page out, with an em-dash --- and
quotation marks ``like these'' and a footnote.\footnote{A footnote, which is
an insertion, so the page builder must split it.}
\subsection{A subsection}
Running text again, with \textit{italics}, \textbf{bold}, and
\texttt{typewriter}, and a formula $\alpha\beta\gamma\delta$ set inline so
that the math list machinery runs inside a paragraph. A displayed formula
follows:
\[ \lim_{n\to\infty}\left(1+\frac{1}{n}\right)^{n}=e . \]
And a little more prose to fill the page out, with an em-dash --- and
quotation marks ``like these'' and a footnote.\footnote{A footnote, which is
an insertion, so the page builder must split it.}
\section{Prose that fills lines and breaks pages}
\label{sec:prose11-3}
Ordinary running text is what a typesetting engine spends most of its effort
upon, so the training document supplies a great deal of it. The words are
chosen for their length and their hyphenation, with incomprehensibility and
antidisestablishmentarianism thrown in so that the patterns are exercised.
Cross references such as \cref{sec:prose11-3} and citations of a page such
as \pageref{sec:prose11-3} make the auxiliary machinery run, and a
paragraph long enough to break over several lines gives the line breaker
something to weigh.
\begin{proposition}\label{prop:11-3}
For every $n\in\mathbb{N}$ the sum $\sum_{k=1}^{n}k$ equals
$\tfrac{1}{2}n(n+1)$.
\end{proposition}
\begin{proof}
Induction on $n$; the base case is $\pairs{1}{1}$ and the step adds $n+1$ to
both sides of the identity, which is what \cref{prop:11-3} asserts.
\end{proof}
\begin{equation}\label{eq:11-3}
  \int_{0}^{\infty}\frac{\vecn{x}^{s-1}}{e^{\vecn{x}}-1}\,d\vecn{x}
    = \Gamma(s)\zeta(s),\qquad \Re(s)>1 .
\end{equation}
\begin{align}
  \mathcal{L}[f](s) &= \int_{0}^{\infty} e^{-st} f(t)\,dt, \\
  \widehat{f}(\xi)  &= \int_{-\infty}^{\infty} f(x) e^{-2\pi i x\xi}\,dx .
\end{align}
\par\noindent\begin{center}
\begin{tabular}{lrrr}
\hline
Name & One & Two & Three \\
\hline
Alpha & 1 & 2 & 3 \\
Beta & 4 & 5 & 6 \\
Gamma & 7 & 8 & 9 \\
\hline
\end{tabular}
\end{center}
\begin{itemize}
  \item A first item with a formula $a^2+b^2=c^2$ in it.
  \item A second item that refers to \eqref{eq:11-3}.
  \item A third item, longer, so that the list breaks across a line and the
    paragraph builder has something to weigh when it chooses the break.
\end{itemize}
\subsection{A subsection}
Running text again, with \textit{italics}, \textbf{bold}, and
\texttt{typewriter}, and a formula $\alpha\beta\gamma\delta$ set inline so
that the math list machinery runs inside a paragraph. A displayed formula
follows:
\[ \lim_{n\to\infty}\left(1+\frac{1}{n}\right)^{n}=e . \]
And a little more prose to fill the page out, with an em-dash --- and
quotation marks ``like these'' and a footnote.\footnote{A footnote, which is
an insertion, so the page builder must split it.}
\subsection{A subsection}
Running text again, with \textit{italics}, \textbf{bold}, and
\texttt{typewriter}, and a formula $\alpha\beta\gamma\delta$ set inline so
that the math list machinery runs inside a paragraph. A displayed formula
follows:
\[ \lim_{n\to\infty}\left(1+\frac{1}{n}\right)^{n}=e . \]
And a little more prose to fill the page out, with an em-dash --- and
quotation marks ``like these'' and a footnote.\footnote{A footnote, which is
an insertion, so the page builder must split it.}
\subsection{A subsection}
Running text again, with \textit{italics}, \textbf{bold}, and
\texttt{typewriter}, and a formula $\alpha\beta\gamma\delta$ set inline so
that the math list machinery runs inside a paragraph. A displayed formula
follows:
\[ \lim_{n\to\infty}\left(1+\frac{1}{n}\right)^{n}=e . \]
And a little more prose to fill the page out, with an em-dash --- and
quotation marks ``like these'' and a footnote.\footnote{A footnote, which is
an insertion, so the page builder must split it.}
\section{Prose that fills lines and breaks pages}
\label{sec:prose11-4}
Ordinary running text is what a typesetting engine spends most of its effort
upon, so the training document supplies a great deal of it. The words are
chosen for their length and their hyphenation, with incomprehensibility and
antidisestablishmentarianism thrown in so that the patterns are exercised.
Cross references such as \cref{sec:prose11-4} and citations of a page such
as \pageref{sec:prose11-4} make the auxiliary machinery run, and a
paragraph long enough to break over several lines gives the line breaker
something to weigh.
\begin{proposition}\label{prop:11-4}
For every $n\in\mathbb{N}$ the sum $\sum_{k=1}^{n}k$ equals
$\tfrac{1}{2}n(n+1)$.
\end{proposition}
\begin{proof}
Induction on $n$; the base case is $\pairs{1}{1}$ and the step adds $n+1$ to
both sides of the identity, which is what \cref{prop:11-4} asserts.
\end{proof}
\begin{equation}\label{eq:11-4}
  \int_{0}^{\infty}\frac{\vecn{x}^{s-1}}{e^{\vecn{x}}-1}\,d\vecn{x}
    = \Gamma(s)\zeta(s),\qquad \Re(s)>1 .
\end{equation}
\begin{align}
  \mathcal{L}[f](s) &= \int_{0}^{\infty} e^{-st} f(t)\,dt, \\
  \widehat{f}(\xi)  &= \int_{-\infty}^{\infty} f(x) e^{-2\pi i x\xi}\,dx .
\end{align}
\par\noindent\begin{center}
\begin{tabular}{lrrr}
\hline
Name & One & Two & Three \\
\hline
Alpha & 1 & 2 & 3 \\
Beta & 4 & 5 & 6 \\
Gamma & 7 & 8 & 9 \\
\hline
\end{tabular}
\end{center}
\begin{itemize}
  \item A first item with a formula $a^2+b^2=c^2$ in it.
  \item A second item that refers to \eqref{eq:11-4}.
  \item A third item, longer, so that the list breaks across a line and the
    paragraph builder has something to weigh when it chooses the break.
\end{itemize}
\subsection{A subsection}
Running text again, with \textit{italics}, \textbf{bold}, and
\texttt{typewriter}, and a formula $\alpha\beta\gamma\delta$ set inline so
that the math list machinery runs inside a paragraph. A displayed formula
follows:
\[ \lim_{n\to\infty}\left(1+\frac{1}{n}\right)^{n}=e . \]
And a little more prose to fill the page out, with an em-dash --- and
quotation marks ``like these'' and a footnote.\footnote{A footnote, which is
an insertion, so the page builder must split it.}
\subsection{A subsection}
Running text again, with \textit{italics}, \textbf{bold}, and
\texttt{typewriter}, and a formula $\alpha\beta\gamma\delta$ set inline so
that the math list machinery runs inside a paragraph. A displayed formula
follows:
\[ \lim_{n\to\infty}\left(1+\frac{1}{n}\right)^{n}=e . \]
And a little more prose to fill the page out, with an em-dash --- and
quotation marks ``like these'' and a footnote.\footnote{A footnote, which is
an insertion, so the page builder must split it.}
\subsection{A subsection}
Running text again, with \textit{italics}, \textbf{bold}, and
\texttt{typewriter}, and a formula $\alpha\beta\gamma\delta$ set inline so
that the math list machinery runs inside a paragraph. A displayed formula
follows:
\[ \lim_{n\to\infty}\left(1+\frac{1}{n}\right)^{n}=e . \]
And a little more prose to fill the page out, with an em-dash --- and
quotation marks ``like these'' and a footnote.\footnote{A footnote, which is
an insertion, so the page builder must split it.}
\chapter{A chapter of ordinary matter, number 12}
\section{Prose that fills lines and breaks pages}
\label{sec:prose12-1}
Ordinary running text is what a typesetting engine spends most of its effort
upon, so the training document supplies a great deal of it. The words are
chosen for their length and their hyphenation, with incomprehensibility and
antidisestablishmentarianism thrown in so that the patterns are exercised.
Cross references such as \cref{sec:prose12-1} and citations of a page such
as \pageref{sec:prose12-1} make the auxiliary machinery run, and a
paragraph long enough to break over several lines gives the line breaker
something to weigh.
\begin{proposition}\label{prop:12-1}
For every $n\in\mathbb{N}$ the sum $\sum_{k=1}^{n}k$ equals
$\tfrac{1}{2}n(n+1)$.
\end{proposition}
\begin{proof}
Induction on $n$; the base case is $\pairs{1}{1}$ and the step adds $n+1$ to
both sides of the identity, which is what \cref{prop:12-1} asserts.
\end{proof}
\begin{equation}\label{eq:12-1}
  \int_{0}^{\infty}\frac{\vecn{x}^{s-1}}{e^{\vecn{x}}-1}\,d\vecn{x}
    = \Gamma(s)\zeta(s),\qquad \Re(s)>1 .
\end{equation}
\begin{align}
  \mathcal{L}[f](s) &= \int_{0}^{\infty} e^{-st} f(t)\,dt, \\
  \widehat{f}(\xi)  &= \int_{-\infty}^{\infty} f(x) e^{-2\pi i x\xi}\,dx .
\end{align}
\par\noindent\begin{center}
\begin{tabular}{lrrr}
\hline
Name & One & Two & Three \\
\hline
Alpha & 1 & 2 & 3 \\
Beta & 4 & 5 & 6 \\
Gamma & 7 & 8 & 9 \\
\hline
\end{tabular}
\end{center}
\begin{itemize}
  \item A first item with a formula $a^2+b^2=c^2$ in it.
  \item A second item that refers to \eqref{eq:12-1}.
  \item A third item, longer, so that the list breaks across a line and the
    paragraph builder has something to weigh when it chooses the break.
\end{itemize}
\subsection{A subsection}
Running text again, with \textit{italics}, \textbf{bold}, and
\texttt{typewriter}, and a formula $\alpha\beta\gamma\delta$ set inline so
that the math list machinery runs inside a paragraph. A displayed formula
follows:
\[ \lim_{n\to\infty}\left(1+\frac{1}{n}\right)^{n}=e . \]
And a little more prose to fill the page out, with an em-dash --- and
quotation marks ``like these'' and a footnote.\footnote{A footnote, which is
an insertion, so the page builder must split it.}
\subsection{A subsection}
Running text again, with \textit{italics}, \textbf{bold}, and
\texttt{typewriter}, and a formula $\alpha\beta\gamma\delta$ set inline so
that the math list machinery runs inside a paragraph. A displayed formula
follows:
\[ \lim_{n\to\infty}\left(1+\frac{1}{n}\right)^{n}=e . \]
And a little more prose to fill the page out, with an em-dash --- and
quotation marks ``like these'' and a footnote.\footnote{A footnote, which is
an insertion, so the page builder must split it.}
\subsection{A subsection}
Running text again, with \textit{italics}, \textbf{bold}, and
\texttt{typewriter}, and a formula $\alpha\beta\gamma\delta$ set inline so
that the math list machinery runs inside a paragraph. A displayed formula
follows:
\[ \lim_{n\to\infty}\left(1+\frac{1}{n}\right)^{n}=e . \]
And a little more prose to fill the page out, with an em-dash --- and
quotation marks ``like these'' and a footnote.\footnote{A footnote, which is
an insertion, so the page builder must split it.}
\section{Prose that fills lines and breaks pages}
\label{sec:prose12-2}
Ordinary running text is what a typesetting engine spends most of its effort
upon, so the training document supplies a great deal of it. The words are
chosen for their length and their hyphenation, with incomprehensibility and
antidisestablishmentarianism thrown in so that the patterns are exercised.
Cross references such as \cref{sec:prose12-2} and citations of a page such
as \pageref{sec:prose12-2} make the auxiliary machinery run, and a
paragraph long enough to break over several lines gives the line breaker
something to weigh.
\begin{proposition}\label{prop:12-2}
For every $n\in\mathbb{N}$ the sum $\sum_{k=1}^{n}k$ equals
$\tfrac{1}{2}n(n+1)$.
\end{proposition}
\begin{proof}
Induction on $n$; the base case is $\pairs{1}{1}$ and the step adds $n+1$ to
both sides of the identity, which is what \cref{prop:12-2} asserts.
\end{proof}
\begin{equation}\label{eq:12-2}
  \int_{0}^{\infty}\frac{\vecn{x}^{s-1}}{e^{\vecn{x}}-1}\,d\vecn{x}
    = \Gamma(s)\zeta(s),\qquad \Re(s)>1 .
\end{equation}
\begin{align}
  \mathcal{L}[f](s) &= \int_{0}^{\infty} e^{-st} f(t)\,dt, \\
  \widehat{f}(\xi)  &= \int_{-\infty}^{\infty} f(x) e^{-2\pi i x\xi}\,dx .
\end{align}
\par\noindent\begin{center}
\begin{tabular}{lrrr}
\hline
Name & One & Two & Three \\
\hline
Alpha & 1 & 2 & 3 \\
Beta & 4 & 5 & 6 \\
Gamma & 7 & 8 & 9 \\
\hline
\end{tabular}
\end{center}
\begin{itemize}
  \item A first item with a formula $a^2+b^2=c^2$ in it.
  \item A second item that refers to \eqref{eq:12-2}.
  \item A third item, longer, so that the list breaks across a line and the
    paragraph builder has something to weigh when it chooses the break.
\end{itemize}
\subsection{A subsection}
Running text again, with \textit{italics}, \textbf{bold}, and
\texttt{typewriter}, and a formula $\alpha\beta\gamma\delta$ set inline so
that the math list machinery runs inside a paragraph. A displayed formula
follows:
\[ \lim_{n\to\infty}\left(1+\frac{1}{n}\right)^{n}=e . \]
And a little more prose to fill the page out, with an em-dash --- and
quotation marks ``like these'' and a footnote.\footnote{A footnote, which is
an insertion, so the page builder must split it.}
\subsection{A subsection}
Running text again, with \textit{italics}, \textbf{bold}, and
\texttt{typewriter}, and a formula $\alpha\beta\gamma\delta$ set inline so
that the math list machinery runs inside a paragraph. A displayed formula
follows:
\[ \lim_{n\to\infty}\left(1+\frac{1}{n}\right)^{n}=e . \]
And a little more prose to fill the page out, with an em-dash --- and
quotation marks ``like these'' and a footnote.\footnote{A footnote, which is
an insertion, so the page builder must split it.}
\subsection{A subsection}
Running text again, with \textit{italics}, \textbf{bold}, and
\texttt{typewriter}, and a formula $\alpha\beta\gamma\delta$ set inline so
that the math list machinery runs inside a paragraph. A displayed formula
follows:
\[ \lim_{n\to\infty}\left(1+\frac{1}{n}\right)^{n}=e . \]
And a little more prose to fill the page out, with an em-dash --- and
quotation marks ``like these'' and a footnote.\footnote{A footnote, which is
an insertion, so the page builder must split it.}
\section{Prose that fills lines and breaks pages}
\label{sec:prose12-3}
Ordinary running text is what a typesetting engine spends most of its effort
upon, so the training document supplies a great deal of it. The words are
chosen for their length and their hyphenation, with incomprehensibility and
antidisestablishmentarianism thrown in so that the patterns are exercised.
Cross references such as \cref{sec:prose12-3} and citations of a page such
as \pageref{sec:prose12-3} make the auxiliary machinery run, and a
paragraph long enough to break over several lines gives the line breaker
something to weigh.
\begin{proposition}\label{prop:12-3}
For every $n\in\mathbb{N}$ the sum $\sum_{k=1}^{n}k$ equals
$\tfrac{1}{2}n(n+1)$.
\end{proposition}
\begin{proof}
Induction on $n$; the base case is $\pairs{1}{1}$ and the step adds $n+1$ to
both sides of the identity, which is what \cref{prop:12-3} asserts.
\end{proof}
\begin{equation}\label{eq:12-3}
  \int_{0}^{\infty}\frac{\vecn{x}^{s-1}}{e^{\vecn{x}}-1}\,d\vecn{x}
    = \Gamma(s)\zeta(s),\qquad \Re(s)>1 .
\end{equation}
\begin{align}
  \mathcal{L}[f](s) &= \int_{0}^{\infty} e^{-st} f(t)\,dt, \\
  \widehat{f}(\xi)  &= \int_{-\infty}^{\infty} f(x) e^{-2\pi i x\xi}\,dx .
\end{align}
\par\noindent\begin{center}
\begin{tabular}{lrrr}
\hline
Name & One & Two & Three \\
\hline
Alpha & 1 & 2 & 3 \\
Beta & 4 & 5 & 6 \\
Gamma & 7 & 8 & 9 \\
\hline
\end{tabular}
\end{center}
\begin{itemize}
  \item A first item with a formula $a^2+b^2=c^2$ in it.
  \item A second item that refers to \eqref{eq:12-3}.
  \item A third item, longer, so that the list breaks across a line and the
    paragraph builder has something to weigh when it chooses the break.
\end{itemize}
\subsection{A subsection}
Running text again, with \textit{italics}, \textbf{bold}, and
\texttt{typewriter}, and a formula $\alpha\beta\gamma\delta$ set inline so
that the math list machinery runs inside a paragraph. A displayed formula
follows:
\[ \lim_{n\to\infty}\left(1+\frac{1}{n}\right)^{n}=e . \]
And a little more prose to fill the page out, with an em-dash --- and
quotation marks ``like these'' and a footnote.\footnote{A footnote, which is
an insertion, so the page builder must split it.}
\subsection{A subsection}
Running text again, with \textit{italics}, \textbf{bold}, and
\texttt{typewriter}, and a formula $\alpha\beta\gamma\delta$ set inline so
that the math list machinery runs inside a paragraph. A displayed formula
follows:
\[ \lim_{n\to\infty}\left(1+\frac{1}{n}\right)^{n}=e . \]
And a little more prose to fill the page out, with an em-dash --- and
quotation marks ``like these'' and a footnote.\footnote{A footnote, which is
an insertion, so the page builder must split it.}
\subsection{A subsection}
Running text again, with \textit{italics}, \textbf{bold}, and
\texttt{typewriter}, and a formula $\alpha\beta\gamma\delta$ set inline so
that the math list machinery runs inside a paragraph. A displayed formula
follows:
\[ \lim_{n\to\infty}\left(1+\frac{1}{n}\right)^{n}=e . \]
And a little more prose to fill the page out, with an em-dash --- and
quotation marks ``like these'' and a footnote.\footnote{A footnote, which is
an insertion, so the page builder must split it.}
\section{Prose that fills lines and breaks pages}
\label{sec:prose12-4}
Ordinary running text is what a typesetting engine spends most of its effort
upon, so the training document supplies a great deal of it. The words are
chosen for their length and their hyphenation, with incomprehensibility and
antidisestablishmentarianism thrown in so that the patterns are exercised.
Cross references such as \cref{sec:prose12-4} and citations of a page such
as \pageref{sec:prose12-4} make the auxiliary machinery run, and a
paragraph long enough to break over several lines gives the line breaker
something to weigh.
\begin{proposition}\label{prop:12-4}
For every $n\in\mathbb{N}$ the sum $\sum_{k=1}^{n}k$ equals
$\tfrac{1}{2}n(n+1)$.
\end{proposition}
\begin{proof}
Induction on $n$; the base case is $\pairs{1}{1}$ and the step adds $n+1$ to
both sides of the identity, which is what \cref{prop:12-4} asserts.
\end{proof}
\begin{equation}\label{eq:12-4}
  \int_{0}^{\infty}\frac{\vecn{x}^{s-1}}{e^{\vecn{x}}-1}\,d\vecn{x}
    = \Gamma(s)\zeta(s),\qquad \Re(s)>1 .
\end{equation}
\begin{align}
  \mathcal{L}[f](s) &= \int_{0}^{\infty} e^{-st} f(t)\,dt, \\
  \widehat{f}(\xi)  &= \int_{-\infty}^{\infty} f(x) e^{-2\pi i x\xi}\,dx .
\end{align}
\par\noindent\begin{center}
\begin{tabular}{lrrr}
\hline
Name & One & Two & Three \\
\hline
Alpha & 1 & 2 & 3 \\
Beta & 4 & 5 & 6 \\
Gamma & 7 & 8 & 9 \\
\hline
\end{tabular}
\end{center}
\begin{itemize}
  \item A first item with a formula $a^2+b^2=c^2$ in it.
  \item A second item that refers to \eqref{eq:12-4}.
  \item A third item, longer, so that the list breaks across a line and the
    paragraph builder has something to weigh when it chooses the break.
\end{itemize}
\subsection{A subsection}
Running text again, with \textit{italics}, \textbf{bold}, and
\texttt{typewriter}, and a formula $\alpha\beta\gamma\delta$ set inline so
that the math list machinery runs inside a paragraph. A displayed formula
follows:
\[ \lim_{n\to\infty}\left(1+\frac{1}{n}\right)^{n}=e . \]
And a little more prose to fill the page out, with an em-dash --- and
quotation marks ``like these'' and a footnote.\footnote{A footnote, which is
an insertion, so the page builder must split it.}
\subsection{A subsection}
Running text again, with \textit{italics}, \textbf{bold}, and
\texttt{typewriter}, and a formula $\alpha\beta\gamma\delta$ set inline so
that the math list machinery runs inside a paragraph. A displayed formula
follows:
\[ \lim_{n\to\infty}\left(1+\frac{1}{n}\right)^{n}=e . \]
And a little more prose to fill the page out, with an em-dash --- and
quotation marks ``like these'' and a footnote.\footnote{A footnote, which is
an insertion, so the page builder must split it.}
\subsection{A subsection}
Running text again, with \textit{italics}, \textbf{bold}, and
\texttt{typewriter}, and a formula $\alpha\beta\gamma\delta$ set inline so
that the math list machinery runs inside a paragraph. A displayed formula
follows:
\[ \lim_{n\to\infty}\left(1+\frac{1}{n}\right)^{n}=e . \]
And a little more prose to fill the page out, with an em-dash --- and
quotation marks ``like these'' and a footnote.\footnote{A footnote, which is
an insertion, so the page builder must split it.}
\end{document}
